\documentclass[12pt,a4paper]{article}
\usepackage[T1]{fontenc}
\usepackage[utf8]{inputenc}
\usepackage[ngerman,latin]{babel}
\usepackage{amsmath,amssymb,amsfonts}
\usepackage{array,booktabs,longtable,multirow}
\usepackage{geometry}
\usepackage{fancyhdr}
\usepackage{titlesec}
\usepackage{etoolbox}
\usepackage{microtype}
\usepackage{enumitem}
\usepackage{graphicx}

\geometry{margin=2.5cm,top=3cm,bottom=3cm}

\pagestyle{fancy}
\fancyhf{}
\fancyhead[LE,RO]{\thepage}
\fancyhead[CE]{\nouppercase{\leftmark}}
\renewcommand{\headrulewidth}{0.4pt}

\titleformat{\section}{\large\bfseries\centering}{\thesection}{1em}{}
\titleformat{\subsection}{\normalsize\bfseries\centering}{\thesubsection}{1em}{}

\newcommand{\gd}{\mathrm{g}}
\newcommand{\std}{\mathrm{st}}
\newcommand{\minn}{\mathrm{m}}
\newcommand{\sek}{\mathrm{s}}
\newcommand{\zd}{\mathrm{Z}.\,\mathrm{D}.}

\title{Carl Friedrich Gauß Werke --- Band 6, Teil 9}
\author{Carl Friedrich Gauß}
\date{}

\begin{document}
\maketitle
\thispagestyle{empty}
\newpage
\setcounter{page}{391}

\section*{COMET.}
\markboth{COMET.}{}
\addcontentsline{toc}{section}{Comet}

Die Vergleichung der letztern Elemente mit den sp"atesten hiesigen Beobachtungen gab folgende
Unterschiede:

\begin{center}
\begin{tabular}{l|c|c}
\toprule
& ger.~Aufst. & Abweichung \\
\midrule
Juni 30 & $-\,50''4$ & $+\,15''9$ \\
Juli 13 & $-\,16.6$ & $-\,44.1$ \\
\phantom{13} 27 & $-\,23.8$ & $-\,47.8$ \\
\phantom{13} 29 & $+\,3.6$ & $-\,47.5$ \\
August \phantom{0}4 & $-\,6.5$ & $-\,38.1$ \\
\phantom{13} 25 & $-\,2.4$ & $+\,15.9$ \\
\bottomrule
\end{tabular}
\end{center}

Wir holen bei dieser Gelegenheit auch noch die Anzeige der Beobachtungen der \textit{Juno} nach,
welche vom Hrn.~Prof.~\textsc{Gauss} um die Zeit der Opposition gemacht sind. Es scheint nicht, dass dieser
diesmal ungemein lichtschwache Planet auf einer andern Sternwarte beobachtet w"are.

\begin{center}
\begin{tabular}{r@{~}l|c|c|c}
\toprule
\multicolumn{2}{c|}{1815} & Mittlere Zeit & Scheinb.~gerade & Scheinbare \\
\multicolumn{2}{c|}{} & in G"ottingen & Aufsteigung & Abweichung \\
\midrule
M"arz & 1  & $10^{\mathrm{h}}\;22^{\mathrm{m}}\;0^{\mathrm{s}}5$ & $197^{\circ}\;7'\;56''7$ & $2^{\circ}\;19'\;3''0$~S. \\
     & 29 & $10\;\phantom{0}8\;57.0$ & $192\;31\;51.8$ & $1\;32\;30.1$~N. \\
April & 8 & $\phantom{0}9\;27\;42.9$ & $190\;34\;51.9$ & $2\;53\;24.2$~N. \\
\bottomrule
\end{tabular}
\end{center}

Hr.~\textsc{Nicolai} "ubernahm die Vergleichung dieser Beobachtungen mit den letzten durch Hrn.~\textsc{M"obius}
berechneten Elementen (S.~Mon.~Corresp.~1813. Bd.~XXVIII. S.~577), und fand, nachdem er die Epoche
um $4'\,55''$ vergr"ossert hatte, folgende

\begin{center}
\begin{tabular}{l|c|c|c|c}
\toprule
& \multicolumn{2}{c|}{Unterschiede} & & \\
& ger.~Aufst. & Abweich. & L"ange & Breite \\
\midrule
M"arz 1  & $-\,20''6$ & $-\,49.7$ & $0$ & $-\,53''7$ \\
\phantom{M"arz} 29 & $-\,12.2$ & $-\,50.5$ & $+\,8.6$ & $-\,51.2$ \\
April 8 & $-\,16.3$ & $-\,45.1$ & $+\,2.9$ & $-\,47.8$ \\
\bottomrule
\end{tabular}
\end{center}

und hieraus die Opposition mit der Sonne

\begin{center}
\begin{tabular}{l}
1815 M"arz 31.\quad $12^{\mathrm{h}}\;56^{\mathrm{m}}\;43^{\mathrm{s}}$ Mittlere Zeit in G"ottingen \\
Wahre L"ange \dotfill $190^{\circ}\;25'\;6''6$ \\
N"ordliche Geocentrische Breite \dotfill $6\;29\;10.0$ \\
\end{tabular}
\end{center}

Ferner erhielt dieser geschickte Astronom durch Verbindung dieser Opposition mit den drei vorhergehenden folgende

\begin{center}
\textbf{Neue Elemente der Juno.}
\end{center}

\begin{center}
\begin{tabular}{l@{\dotfill}l}
Epoche der mittlern L"ange 1816 & $230^{\circ}\;11'\;34''2$ f"ur den Meridian von G"ottingen \\
T"agliche tropische Bewegung & $812''9304$ \\
L"ange des Perihels 1816 & $53^{\circ}\;14'\;53''8$ \\
L"ange des aufsteigenden Knoten 1816 & $171\;\phantom{0}9\;58.9$ \\
Neigung der Bahn & $13\;14\;0.1$ \\
Excentricit"atswinkel & $14\;43\;28.84$ \\
Logarithm.~der halben grossen Axe & $0.4266844$ \\
\end{tabular}
\end{center}

Eine nach diesen Elementen von Hrn~\textsc{Nicolai} berechnete Ephemeride wird in dem astronomischen
Jahrbuche f"ur 1818 bekannt gemacht werden.

\newpage

\section*{JUNO.}
\markboth{JUNO.}{}

\subsection*{GAUSS AN BODE.}
\addcontentsline{toc}{subsection}{GAUSS AN BODE.}

\begin{center}
\rule{0.6\textwidth}{0.4pt}\\[0.5em]
\textit{Astronomisches Jahrbuch f"ur 1819. Seite 219\,..\,221. Berlin 1816.}\\[0.3em]
\rule{0.6\textwidth}{0.4pt}
\end{center}

\begin{flushright}
G"ottingen den 10.~August 1816.
\end{flushright}

Ihrem Wunsche zufolge, habe ich das Vergn"ugen, Ihnen, verehrterester Freund, beigehend Ephemeriden
f"ur die n"achste Erscheinung der Pallas und Juno zu "ubersenden, die ich mit einigen Anmerkungen begleite.

Die Ephemeride f"ur die Pallas ist von Hrn~\textsc{Westphal} aus Schwerin berechnet, welcher sich hier
gegenw"artig mit Eifer der Astronomie widmet. Es ist dabei --- wie immer seit 5 Jahren --- auf die
St"orungen nach meiner Theorie R"ucksicht genommen, und daher die genaueste Uebereinstimmung zu erwarten.
Die n"achste Opposition hat Hr.~\textsc{Westphal}, wie folgt, im Voraus berechnet:

\begin{center}
\begin{tabular}{l@{\quad}l}
1817 Juli 11. & $8^{\mathrm{h}}\;45^{\mathrm{m}}\;5^{\mathrm{s}}$ Mittlere Zeit in G"ottingen \\
Wahre L"ange \dotfill & $289^{\circ}\;3'\;22''$ \\
Geocentrische Breite \dotfill & $43\;33\;\phantom{0}1$ \\
\multicolumn{2}{c}{$\rho = 0.05997$}
\end{tabular}
\end{center}

An Licht wird die Pallas im n"achsten Jahre sehr schwach sein, und nur 10--11.~Gr"osse haben.
Folgendes ist die Uebersicht der Lichtst"arke in den letzten und der k"unftigen Opposition:

\begin{center}
\begin{tabular}{c|c|c|c}
1812 & 0.01666 & 1816 & 0.05997 \\
1813 & 0.01475 & 1817 & 0.01289 \\
1814 & 0.05476 & 1818 & 0.02122 \\
\end{tabular}
\end{center}

Die Juno habe ich selbst zweimal am Kreismikrometer beobachtet:

\begin{center}
\begin{tabular}{c|c|c|c}
1816 & Mittlere Zeit & Gerade Aufsteig. & S"udl.~Abweich. \\
\midrule
Juni 12 & $10^{\mathrm{h}}\;47^{\mathrm{m}}\;27^{\mathrm{s}}5$ & $251^{\circ}\;46'\;8''0$ & $3^{\circ}\;42'\;4''7$ \\
\phantom{12} 13 & $10\;27\;34.7$ & $251\;33\;47.6$ & $3\;41\;11.7$ \\
\end{tabular}
\end{center}

Bei dem so ungew"ohnlich schlechten Wetter, welches im vorigen Fr"uhjahr herrschte, konnten nicht mehr
Beobachtungen gemacht werden. Ich habe daraus folgende Resultate f"ur die Opposition abgeleitet:

\begin{center}
\begin{tabular}{l@{\quad}l}
1816 Juni 3. & $4^{\mathrm{h}}\;20^{\mathrm{m}}\;20^{\mathrm{s}}$ Mittlere Zeit in G"ottingen \\
Wahre L"ange \dotfill & $252^{\circ}\;51'\;39''2$ \\
Geocentrische Breite \dotfill & $18\;34\;3.1$~N"ordl.
\end{tabular}
\end{center}

Hr.~\textsc{Encke}, welcher nachher diese Rechnung wiederholt und dabei auch die Wilnaer Beobachtungen
mit benutzt hat, fand

\begin{center}
\begin{tabular}{l@{\quad}l}
Juni 3 & $4^{\mathrm{h}}\;17^{\mathrm{m}}\;31^{\mathrm{s}}$ Mittlere Zeit in G"ottingen \\
& $252^{\circ}\;51'\;32''6$ \\
& $18\;34\;14.6$
\end{tabular}
\end{center}

Aus der Combination der vier letzten Oppositionen leitete ich folgende Elemente ab:

\newpage

\begin{center}
\begin{tabular}{l@{\dotfill}l}
Mittlere L"ange 1817 Sept.~6.\ $0^{\mathrm{h}}$ in G"ottingen & $8^{\circ}\;51'\;41''$ \\
Sonnenn"ahe & $53\;16\;46$ \\
Knoten & $171\;\phantom{0}9\;41$ \\
Neigung & $13\;\phantom{0}3\;49$ \\
Excentricit"atswinkel & $14\;46\;\phantom{0}0$ \\
T"agliche mittlere tropische Bewegung & $812''40187$ \\
Logarithm der halben grossen Axe & $0.4269223$ \\
\end{tabular}
\end{center}

(Es ist wohl "uberfl"ussig zu bemerken, dass diese Elemente blos zu dem Zweck berechnet sind
zu der Berechnung der n"achsten Ephemeride zu dienen: so lange die St"orungen noch nicht entwickelt
sind, ist dies Verfahren das bequemste und genaueste.)

Nach diesen Elementen hat Hr.~Prof.~\textsc{Harding} eine Ephemeride f"ur die n"achste Erscheinung berechnet,
von welcher ich Ihnen hier eine Abschrift beilege. Dieselbe Rechnung ist auch von Hrn.~\textsc{Posselt}
ausgef"uhrt, dessen Ephemeride von der gegenw"artigen "uberall nur in den Secunden abweichend
in Hrn.~von~\textsc{Lindenau}'s Zeitschrift abgedruckt wird. Die n"achste Opposition f"allt nach Hrn.~Professor
\textsc{Harding}'s Rechnung aus obigen Elementen:

\begin{center}
\begin{tabular}{l@{\quad}l}
1817 Sept.~6. & $4^{\mathrm{h}}\;10^{\mathrm{m}}\;30^{\mathrm{s}}$ mittlere Zeit in G"ottingen. \\
Wahre L"ange \dotfill & $343^{\circ}\;37'\;36''$ \\
Geocentrische Breite \dotfill & $3\;\phantom{0}6\;\phantom{0}9$~N"ordl. \\
Die Lichtst"arke 1816 \dotfill & $0.01747$ \\
\phantom{Die Lichtst"arke} 1817 \dotfill & $0.11398$ \\
\end{tabular}
\end{center}

Erlauben Sie mir noch ein Paar Bemerkungen "uber ver"anderliche Sterne. Dem von Herrn Prof.~\textsc{Harding}
gefundenen in der Jungfrau legt Hr.~\textsc{Koch} im Jahrbuche 1818 eine Periode von 295 Tagen bei
ohne anzugeben, worauf sich dieselbe gr"undet: allein mit den Beobachtungen des Hrn.~Prof.~\textsc{Harding}
vertr"agt sich diese Periode nicht. --- Einen zweiten ver"anderlichen Stern hat Hr.~Prof.~\textsc{Harding} im
Jahre 1811 im Wassermann s"udlich unter $\tau$ bemerkt, der von der sechsten Gr"osse allm"ahlich, und zwar
in ungef"ahr $2\frac{1}{2}$ Monaten zur v"olligen Unsichtbarkeit "ubergeht, dessen Lichtwechsel "uberhaupt aber sich
in keine regelm"assige Periode zu f"ugen scheint. Seine mittlere Stellung 1800 war

\begin{center}
\begin{tabular}{l@{\dotfill}l}
Gerade Aufsteigung & $353^{\circ}\;21'\;48''$ \\
J"ahrliche Ver"anderung & $+3^{\mathrm{s}}\;66$ \\
S"udliche Abweichung & $16^{\circ}\;23'\;13''4$ \\
J"ahrliche Abnahme & $19.88$
\end{tabular}
\end{center}

\begin{flushright}
\textsc{V.~Lindenau u.~Bohnenberger} Zeitschrift f"ur Astronomie u.~verwandte Wiss.\\
Band IV. Seite 119\,..\,131. Juli u.~August. 1817.
\end{flushright}

\newpage

\section*{POLARSTERN.}
\markboth{POLARSTERN.}{}
\addcontentsline{toc}{section}{Polarstern}

\begin{flushright}
G"ottingen, den 18.~Juli 1817.
\end{flushright}

Wenn ich nicht irre, habe ich Ihnen, theuerster Freund, meine Beobachtungen des Polarsterns
am zw"olfz"olligen \textsc{Reichenbach}schen Kreise noch nicht vollst"andig mitgetheilt: ich freue mich, dass ich
jetzt eine Reihe Solstitialbeobachtungen damit verbinden kann, die in diesem Jahre durch das sch"onste
Wetter mehr als je beg"unstigt wurden. Diese Beobachtungen sind auf der neuen Sternwarte gemacht,
deren Bau "ubrigens, wie Sie wissen, noch nicht ganz vollendet ist, und deren eigentliche Th"atigkeit erst
von der Zeit datiren wird, wo die neuen Meridian-Instrumente angelangt und aufgestellt sein werden.
Die Mitte der neuen Sternwarte liegt nach meiner geod"atischen Bestimmung $6''11$ s"udlich, und $1^{\mathrm{s}}90$
"ostlich von der Mitte der alten. Der Platz, wo ich den Polarstern beobachtete, liegt $0''12$ n"ordlich, der
Platz der Sonnenbeobachtungen $0''16$ s"udlich von der Mitte.

Sie erinnern sich, dass in einigen meiner fr"ubern Briefe von der Muthmassung die Rede gewesen
ist, der so viel besprochene Unterschied der Polh"ohe aus Beobachtungen der Circumpolarsterne und aus
Sonnenbeobachtungen k"onne bei den kleinern Vervielf"altigungskreisen seinen Grund in einer Biegung
der Objectivseite des Fernrohrs haben, da man am Objectiv als Gegengewicht nicht blos f"ur das
prismatische Ocular, sondern auch f"ur die Klammer einen starken messingenen Ring vorsteckt. Falls
dadurch eine nachtheilige Biegung hervorgebracht werden sollte, so m"usste nach \textsc{Reichenbach}'s Rathe dies
vorgesteckte Gewicht so viel vermindert werden, dass dadurch blos das Fernrohr geh"orig "aquilibrirt
w"urde und die Klammer durch ein am Alhidadenkreise angeschraubtes Gewicht zu balanciren. Zu
meiner Absicht war es jedoch noch zweckm"assiger, jenes vorgesteckte Gewicht ganz wegzunehmen, und
das ganze Gleichgewicht durch ein am Alhidadenkreise angeschraubtes Gewicht zu bewirken. Denn es
ist klar, dass, wenn jene Vermuthung gegr"undet gewesen w"are, nunmehr durch die nicht aufgehobene
Biegung der Ocularseite des Fernrohrs ein Fehler im entgegengesetzten Sinn hervorgehen m"usste. Wenn
vorher wegen einer Biegung alle Zenithdistanzen zu klein gemessen waren, so m"ussten sie nun alle zu
gross werden und Polarstern-Beobachtungen eine zu kleine, die Sonne eine zu grosse Polh"ohe geben.

Allein der Erfolg hat dies durchaus nicht best"atigt; die auf der neuen Sternwarte aus dem Polarstern,
bei angeschraubtem Gewicht, bestimmte Polh"ohe ist fast haarscharf dieselbe, die aus der Uebertragung
der mit angestecktem Gewicht auf der alten Sternwarte beobachteten folgt, und die Sonnenbeobachtungen
weichen in demselben Sinn und um dieselbe Gr"osse ab wie vorher. Ich hatte anfangs, ehe das
Wetter zusammenh"angende astronomische Beobachtungen erlaubte, einen andern Weg gew"ahlt, jene
Muthmassung zu pr"ufen. Ich hatte n"amlich mehrere Reihen Zenithdistanzen eines terrestrischen Gegenstandes
beobachtet, einige mit vorgestecktem, die andern mit angeschraubtem Gewicht. Die ersten
Reihen gaben wirklich einen Unterschied in dem Sinn und ungef"ahr von der Gr"osse wie er sein m"usste,
wenn jene Quelle die wahre gewesen w"are. Allein bei den nachfolgenden Reihen verschwand dieser
Unterschied fast ganz, und "uberhaupt gaben diese Beobachtungen des irdischen Gegenstandes, auf welchen
sich nicht so scharf pointiren liess wie auf Objecte am Himmel, keine so gute Harmonie unter sich,
wie zu dieser Pr"ufung n"othig gewesen w"are.

Ist dieser "ausserst merkw"urdige Unterschied dem Instrument zuzuschreiben, so muss in seinem
Bau etwas sein, was bewirkt, dass wir die Zenithdistanzen damit zu klein beobachten. Jenes Gewicht
war nicht schuld; mehrere andere Umst"ande, die ich in Verdacht zu ziehen geneigt war, habe ich
gleichfalls nach sorgf"altiger Pr"ufung einflusslos gefunden. Ich gestehe, dass ich bis jetzt keinen gefunden
habe, dem ich eine solche Wirkung beimessen k"onnte. Eine Wirkung entgegengesetzter Art w"urde
viel leichter zu erkl"aren sein, und kann allerdings, wenn man nicht auf seiner Hut ist, stattfinden.
Ich selbst habe, wie Sie wissen, im vorigen Winter eine solche Erfahrung gemacht. Der Kreis war
von mir ganz zerlegt, gereinigt und wieder zusammengesetzt worden. Meine ersten Beobachtungen des
Polarsterns gaben mir wiederholt eine um $20''$ zu kleine Polh"ohe. Es war also etwas da, was die
Zenithdistanzen vermehrte. Ich fand nach mehreren anderen Pr"ufungen die Quelle des Fehlers bald in
einer Biegung der Ocularr"ohre, die an ihrem "aussersten Ende im Fernrohr durch die Feder fest gedr"uckt
wird. Diese mochte etwas lahm geworden sein; nachdem ihr durch einige Biegung die n"othige Spannkraft
wiedergegeben war, fiel der Fehler sogleich weg. Auch wenn der Faden, an dem die Zenithdistanzen
beobachtet werden, nicht geh"orig gespannt ist, wird derselbe Fehler erfolgen. Der Faden h"angt dann
in der Mitte durch sein eignes Gewicht nieder; und man beobachtet offenbar die Zenithdistanzen zu gross.
Es braucht wohl nicht erinnert zu werden, dass hier nicht von groben Fehlern die Rede ist; allein wenn
der schlaffe Faden in der Mitte ein Paar Secunden niederh"angt, w"ahrend seine ganze L"ange vielleicht
"uber einen Grad betr"agt, so erkennt das Auge noch keine Kr"ummung, es sei denn, dass man das Bild
einer vollkommen geraden horizontalen Linie aus einer schicklichen Entfernung sich auf den Faden
projiciren l"asst.

Es ist aber gewiss, dass man das Instrument nicht fehlerfrei sprechen kann, ohne den entgegengesetzten
Fehler andern Instrumenten aufzub"urden, die einen solchen Unterschied zwischen den Stern- und
Sonnenbeobachtungen gar nicht, oder von geringerer Gr"osse geben. Hat der 12 zollige Kreis Recht,
so geben die dreif"ussigen Kreise mit stehender S"aule und \textsc{Bessel}'s Kreis die Zenithdistanzen zu gross;
hat der dreif"ussige Kreis Recht, so gibt der zw"olfzollige die Zenithdistanzen zu klein und der \textsc{Bessel}sche
zu gross; hat der letzte Recht, so geben alle \textsc{Reichenbach}schen Kreise zu kleine Zenithdistanzen.
Allein mit Gewissheit kann man keinem Instrumente Recht geben, wenn man nicht entweder die m"ogliche
Fehlerquelle an den andern Instrumenten bestimmt und unl"augbar nachgewiesen, oder auch anderswoher
die Unrichtigkeit der mit den andern Instrumenten gefundenen Resultate entschieden dargelegt hat.
Wie die Sachen jetzt liegen, m"ussen wir, d"aucht mir, noch gestehen, dass kein bestimmter Ausspruch
zu thun ist. Die relative Wahrscheinlichkeit wird jeder Astronom nach subjectiven Gr"unden vielleicht
anders abw"agen, denn ich finde es allerdings nat"urlich, dass wenn gleich gegen alle Instrumente kein
bestimmter Vorwurf gerichtet werden kann, man dies am lebendigsten in Beziehung auf dasjenige f"uhlt,
womit man selbst beobachtet und womit man also am meisten vertraut ist.

Zwei Bemerkungen aber kann ich doch nicht umhin, hier noch beizuf"ugen.

Erstlich scheint mir, dass man das Zeugniss der zw"olfz"olligen Kreise nicht deswegen unbedingt
verwerfen kann, weil sie die kleinsten von diesen Instrumenten sind, wenn man nicht nachweisen kann,
dass sie wegen ihrer Kleinheit constante Fehler, die immer in demselben Sinn wirken, hervorbringen.
Zuf"allige Fehler, die bald so bald anders wirken, k"onnte man eher wegen der Kleinheit erwarten, und
doch zeigt die Erfahrung, dass sie eine eben so gute, oder wenigstens nicht viel schlechtere Harmonie
geben, wie die gr"ossern Instrumente. Ich weiss nicht, ob man im Allgemeinen nicht bei gr"ossern
Instrumenten wegen ihrer Gr"osse wenigstens eben so sehr auf seiner Hut vor constanten Fehlern sein muss.

Zweitens gestehe ich, dass, ohne alle R"ucksicht auf andere Umst"ande, ein Instrument deswegen,
weil es gar keinen Unterschied zwischen Stern- und Sonnenbeobachtungen gibt, mir noch keine Pr"asumption
f"ur sich zu haben scheint. Nicht zu gedenken, dass es doch noch nicht ganz ausgemacht ist,
ob die Refraction f"ur die Sonne und f"ur die Fixsterne genau dieselbe ist, so kommt hier noch ein
anderer wichtiger Umstand in Betracht, wor"uber ich mit Ihnen, wie Sie sich erinnern, schon vor mehreren
Jahren mich unterhalten habe, und den ich als eine m"ogliche Erkl"arungsart des R"athsels damals
aufsstellte, ohne eben sehr viel Gewicht darauf legen zu wollen. Ich wiederhole hier diesen Gegenstand,
da bisher, so viel ich weiss, noch nicht "offentlich die Rede davon gewesen ist. Wir beobachten den
Mittelpunkt der Figur der Sonne; die Anziehungskraft der Sonne hingegen wird auf die Planeten
proxime so wirken, als w"are sie im Schwerpunkt vereinigt, und der Schwerpunkt der Sonne, nicht der
Mittelpunkt der Figur wird in der Ebene der Erdbahn liegen. Man darf nicht annehmen, dass beide
in mathematischer Sch"arfe zusammenfallen, da ohne Zweifel der Sonnenk"orper keine ganz gleichartige
Masse ist. Fiele der Schwerpunkt etwas s"udlich vom Mittelpunkt der Figur, so w"urde letzterer immer
eine kleine n"ordliche Breite haben, und in so fern man diese ignorirt, m"usste aus Sonnenbeobachtungen
auf der n"ordlichen Erdhalbkugel immer eine zu kleine Polh"ohe folgen.

\newpage

So leicht und nat"urlich inzwischen diese Erkl"arungsart scheint, so ist sie doch noch mit ihren
Schwierigkeiten verbunden. Man m"usste erst in eine tiefere Untersuchung eingehen, in wie fern eine
solche Zusammensetzung des Sonnenk"orpers m"oglich ist, wobei der Schwerpunkt um $\frac{1}{20}$ oder um
$\frac{1}{100}$ des Sonnendurchmessers s"udlich vom Mittelpunkt der Figur fiele, ohne dass die Oberfl"ache des
Sonnenk"orpers, nach den Bedingungen des Gleichgewichts, mehr von der Kugelgestalt abwiche, als es
nach den Beobachtungen der Fall ist. Auch die beobachteten \textsc{Mercurs}-Durchg"ange kommen hier in
Betracht. Da die aus denselben abgeleiteten Breiten des \textsc{Mercurs} sich auf die Voraussetzung gr"unden,
dass der Mittelpunkt der Figur der Sonne in der Ebene der Ecliptik selbst ist, so w"urden alle
beobachteten Breiten zu s"udlich sein. Die daraus abgeleitete L"ange des aufsteigenden Knotens w"urde zu gross,
die des niedersteigenden zu klein ausfallen, und beide w"urden nicht "ubereinstimmen. Die von Ihnen in der
\textsc{Mercur}stheorie mitgetheilten Zahlen deuten in der That auf so etwas hin, scheinen sich aber nicht mit einem
so grossen Unterschiede, wie die zw"olfz"olligen Kreise geben, zu vertragen. Freilich sind die beobachteten
Durchg"ange durch den niedersteigenden Knoten in zu geringer Anzahl. Und da Sie "uber die von Ihnen
gew"ahlten Beobachtungen und die Art, wie die Breiten daraus gefolgert sind, kein Detail mitgetheilt haben,
so l"asst sich "uber den Grad der Genauigkeit, welchen man jeder einzelnen Breite beilegen muss, nicht
urtheilen. Es w"are aber wohl der M"uhe werth, die \textsc{Mercurs}-Durchg"ange aus diesem Gesichtspunkte
noch einmal auf das sorgf"altigste zu discutiren.

Es w"are ungemein interessant, wenn wir Beobachtungen der Sonne und der s"udlichen Circumpolarsterne
mit Repetitionskreisen aus der s"udlichen Erdhalbkugel h"atten. Die s"udliche Breite m"usste aber so gross
sein, dass die Refraction bei den Circumpolarsternen keine erhebliche Ungewissheit hervorbringen k"onnte.

Ich bemerke noch bei dieser Gelegenheit, dass Hr.~Prof.~\textsc{Schumacher} mir die von seinem Bruder,
Hrn.~Capitain \textsc{Schumacher} in Friedrichsmark mit einem \textsc{Reichenbach}schen, dem Ihrigen ganz "ahnlichen,
astronomischen Theodolithen beobachtete Polh"ohe mitgetheilt hat, welche aus Beobachtungen des
Polarsterns $= 55^{\circ}\;58'\;41''3$, aus Sonnenbeobachtungen $= 55^{\circ}\;58'\;36''6$ folgte. Der Unterschied ist fast
genau derselbe wie bei meinen Beobachtungen.

\vspace{1em}

\textbf{1817 Beobachtungen des Polarsterns in der untern Culmination auf der neuen G"ottinger Sternwarte.}

\begin{center}
\begin{tabular}{c|c|c|c|c|c}
\toprule
& & Anz.~d. & Scheinbare & & \\
& & Beob. & Zenithdistanz & Barometer & Thermometer \\
\midrule
M"arz 28 & & 16 & $40^{\circ}\;7'\;22''40$ & $27''\;8.2$ & $+1.4$ \\
April 21 & & 22 & \phantom{0}\phantom{0}\phantom{0}\phantom{0}\phantom{0}\phantom{0}24.98 & & $13.3$ \\
\phantom{21} 2 & & 12 & \phantom{0}\phantom{0}\phantom{0}\phantom{0}\phantom{0}\phantom{0}26.36 & & \\
\phantom{21} 8 & & 28 & \phantom{0}\phantom{0}\phantom{0}\phantom{0}\phantom{0}\phantom{0}23.38 & $10.9$ & $+3.7$ \\
\phantom{21} 11 & & 26 & \phantom{0}\phantom{0}\phantom{0}\phantom{0}\phantom{0}\phantom{0}29.83 & $9.3$ & $+4.4$ \\
Mai \phantom{0}7 & & 12 & \phantom{0}\phantom{0}\phantom{0}\phantom{0}\phantom{0}\phantom{0}25.00 & & $+3.0$ \\
\bottomrule
\end{tabular}
\end{center}

\begin{center}
\textbf{Hieraus folgt:}
\end{center}

\begin{center}
\begin{tabular}{c|c|c|c|c}
\toprule
& & wahre Zenithdist. & Scheinb.~Polardist. & Polh"ohe \\
\midrule
M"arz 28 & & $40^{\circ}\;8'\;13''18$ & $1^{\circ}\;40'\;2''51$ & $51^{\circ}\;31'\;49''33$ \\
April 21 & & $\phantom{40^{\circ}\;}13''87$ & $\phantom{1^{\circ}\;40'\;}\phantom{00}$ & $\phantom{51^{\circ}\;31'\;}49''22$ \\
\phantom{21} 2 & & $\phantom{40^{\circ}\;}15''86$ & $\phantom{1^{\circ}\;40'\;}5''79$ & $\phantom{51^{\circ}\;31'\;}50''13$ \\
\phantom{21} 8 & & $\phantom{40^{\circ}\;}15''76$ & & $\phantom{51^{\circ}\;31'\;}50''03$ \\
\phantom{21} 11 & & $\phantom{40^{\circ}\;}16''51$ & $\phantom{1^{\circ}\;40'\;}3''70$ & $\phantom{51^{\circ}\;31'\;}47''84$ \\
\phantom{21} & & $\phantom{40^{\circ}\;}20''12$ & & $\phantom{51^{\circ}\;31'\;}50''51$ \\
Mai \phantom{0}7 & & $\phantom{40^{\circ}\;}23''90$ & $\phantom{1^{\circ}\;40'\;}4''00$ & $\phantom{51^{\circ}\;31'\;}49''67$ \\
\bottomrule
\end{tabular}
\end{center}

\begin{center}
Also im Mittel aus 158 Beobachtungen $51^{\circ}\;31'\;49''30$, oder auf die Mitte der Sternwarte reducirt
\[51^{\circ}\;31'\;49''30 - 0.644 \cdot Q \quad \text{[Handschriftliche Aufzeichnung.]}\]
\end{center}

Die Refraction ist aus \textsc{Carlini}'s, die Polardistanz aus \textsc{Bessel}'s Tafeln entlehnt. Meine Beobachtungen
des Polarsterns der untern Culmination von 1813 auf der alten Sternwarte, nach denselben Tafeln berechnet,
geben mir die Polh"ohe derselben $51^{\circ}\;31'\;55''48$, also auf die neue Sternwarte reducirt,
136 Beobachtungen,
\[51^{\circ}\;31'\;49''37.\]

Die Ver"anderung des Balancirgewichts hat also gar keinen Einfluss auf die Beobachtungen gehabt.

\newpage

Die obere Culmination des Polarsterns habe ich noch nicht beobachten k"onnen. Am 19.~Februar
beobachtete ich 6 Zenith-Distanzen, 2 Stunden von der obern Culmination entfernt, deren Berechnung
ich hier als ein Beispiel noch beif"ugen will, um mein Versprechen in der M.~C.~1813 S.~484 zu erf"ullen.
Ich sehe mit Vergn"ugen aus dem letzten St"uck Ihrer Zeitschrift, dass auch Hr.~\textsc{Littbow} diese
Beobachtungsart mit Vortheil angewandt hat.

Meine Beobachtungszeiten nach der Uhr waren
\begin{center}
\begin{tabular}{c|c}
$2^{\mathrm{h}}\;54^{\mathrm{m}}\;19^{\mathrm{s}}5$ & $4^{\mathrm{h}}\;24^{\mathrm{m}}\;55^{\mathrm{s}}0$ \\
$3\;\phantom{0}5\;\phantom{0}8.0$ & $3\;32\;57.0$ \\
$3\;17\;36.0$ & $3\;21\;36.0$ \\
\end{tabular}
\end{center}
(Die grossen Intervalle waren Folge des ung"unstigen Wetters, das im ganzen Winter nur in einzelnen
Augenblicken Sterne sichtbar werden liess). Der durchlaufene Bogen war $222^{\circ}\;23'\;53''$, Barometer
$27^{\mathrm{h}}\;10'''3$, Thermometer $+4^{\circ}5$, Culminationszeit des Sterns $0^{\mathrm{h}}\;55^{\mathrm{m}}\;23^{\mathrm{s}}59$.
Mein Berechnungsverfahren will ich hier blos practisch erkl"aren.

Das Mittel aller Uhrzeiten ist $3^{\mathrm{h}}\;6^{\mathrm{m}}\;7^{\mathrm{s}}08$; die Unterschiede von demselben ohne R"ucksicht auf das Zeichen
\begin{center}
\begin{tabular}{c|c}
$11^{\mathrm{m}}\;47^{\mathrm{s}}58$ & $1^{\mathrm{m}}\;10^{\mathrm{s}}08$ \\
$\phantom{0}6\;\phantom{0}12.08$ & $2\;48.92$ \\
$\phantom{0}7\;16.92$ & \\
\end{tabular}
\end{center}

Mit diesen Unterschieden gehe ich in die bekannte Correctionstafel der Circummeridianh"ohen ein
(ich bediene mich der des Hr.~v.~\textsc{Zach} in der Attraction des montagnes), woraus ich erhalte
\begin{center}
\begin{tabular}{c}
$273''01$ \\
$\phantom{0}75.50$ \\
$104.11$ \\
$\phantom{0}\phantom{0}2.68$ \\
$\phantom{0}15.56$ \\
$173.41$ \\
\end{tabular}
\end{center}
das Mittel hieraus ist $107''38$, welchem in derselben Tafel das Argument $7^{\mathrm{m}}\;23^{\mathrm{s}}73$ entspricht.
Ich schliesse hieraus (verm"oge eines kleinen Kunstgriffs, dessen Grund man sich leicht selbst suppliren
kann oder in den G"ott.~gel.~Anz.~1815 S.~454 [M"arz 23 S.~bei den weiter unten folgenden Anzeigen in d.~B.]
nachsehen mag), dass das Mittel der berechneten Zenithdistanzen f"ur die sechs wirklichen Beobachtungszeiten
proxime gleich sei dem Mittel der f"ur die zwei Uhrzeiten
\begin{align*}
&3^{\mathrm{h}}\;6^{\mathrm{m}}\;7^{\mathrm{s}}08 - 7^{\mathrm{m}}\;23^{\mathrm{s}}73 = 2^{\mathrm{h}}\;58^{\mathrm{m}}\;43^{\mathrm{s}}35 \\
&3\;6\;7.08 + 7\;23.73 = 3\;13\;30.81
\end{align*}
welchen die Stundenwinkel $30^{\circ}\;49'\;56''9$ und $34^{\circ}\;31'\;48''3$ entsprechen. Die f"ur diese Stundenwinkel
mit der aus \textsc{Bessel}'s Tafel entlehnten Declination $88^{\circ}\;20'\;7''53$ und der Polh"ohe
$51^{\circ}\;31'\;47''0$ nach bekannten Formeln berechneten Zenithdistanzen sind
\begin{align*}
&37^{\circ}\;2'\;57''31 \\
&37\;\phantom{0}6\;32.75
\end{align*}
also das Mittel $37^{\circ}\;4'\;45''03$ oder mit R"ucksicht auf die durch $d\varphi$, $d\delta$ zu bezeichnenden Correctionen
der zum Grunde gelegten Declination und Polh"ohe{}\footnote{Die Coefficienten sind mit dem mittlern Stundenwinkel aus einer von Hrn.~\textsc{Westphal} f"ur die hiesige Polh"ohe berechneten Tafel entlehnt.}
\[= 37^{\circ}\;4'\;45''03 - 1.00 \cdot d\varphi + 0.83 \cdot d\delta.\]

Das Mittel der beobachteten Zenithdistanz ist $37^{\circ}\;3'\;58''83$, oder von der Refraction $44''64$ befreiet
\[= 37^{\circ}\;4'\;43''47.\]

Hieraus folgt also $d\varphi = +1''56 + 0.83 \cdot d\delta$ oder die Polh"ohe $= 51^{\circ}\;31'\;48''56 + 0.83 \cdot d\delta$.
Die Beobachtungen der untern Culmination haben $51^{\circ}\;31'\;49''42 - d\delta$ gegeben, woraus
$d\delta = +0''47$ folgen w"urde, und die Polh"ohe $51^{\circ}\;31'\;48''95$. (Nat"urlich ist dies nur Beispiels
halber hinzugef"ugt; um die Polh"ohe von der zum Grunde gelegten Declination unabh"angig zu machen,
m"ussten die Beobachtungen in der N"ahe der obern Culmination viel zahlreicher sein, und die Resultate
w"urden dann nach der Methode der kleinsten Quadrate combinirt werden m"ussen). Man k"onnte auch
eine Correction der zum Grunde gelegten geraden Aufsteigung mit einf"uhren, allein man thut besser dies
zu unterlassen, da dies Element gegenw"artig, Dank sei es Ihren und unsers \textsc{Bessel}'s Bem"uhungen,
jetzt so scharf bestimmt ist.

\newpage

\section*{SONNENBEOBACHTUNGEN.}
\markboth{SONNENBEOBACHTUNGEN.}{}
\addcontentsline{toc}{section}{Sonnenbeobachtungen}

Ich lasse nunmehr meine Sonnenbeobachtungen folgen.

\begin{center}
\textbf{1817 Beobachtungen des Sonnenstillstandes.}
\end{center}

\begin{center}
\small
\begin{tabular}{c|c|c|c|c|c|c}
\toprule
& & durchlaufener & Summe der & Anz.~d. & & \\
& & Bogen & Reduction & Beob. & Barometer & Thermometer \\
\midrule
Juni 13 & & & $11'\;28''90$ & & & $19.3$ \\
\phantom{13} 16 & & $283^{\circ}\;11'$ & $\phantom{0}9.75$ & 16 & & $14.3$ \\
\phantom{13} 18 & 23451 & $\phantom{0}21\;\phantom{0}\phantom{0}$ & $38.75$ & 52 & $27^{\mathrm{h}}\;7'''06$ & $20.8$ \\
& 21 & & & & $11.\phantom{0}$ & $6.\phantom{0}$ \\
\phantom{13} 20 & 563 & $2\;46$ & $30.75$ & & $7.7$ & $13.0$ \\
& 2 & & $12\;31.98$ & 16 & $13.1$ & $20.7$ \\
\phantom{13} 24 & 449 & $59\;10.75$ & & & $7.5\;8.9$ & $22.$ \\
& 337 & & $6177\;30.\phantom{00}$ & & & \\
& 562 & $29\;188$ & $25\;45.01$ & 12 & & $20.45$ \\
\phantom{13} 21 & & & $\phantom{0}1.50$ & & $7.4$ & $21.0$ \\
\phantom{13} 27 & & & $19\;20.42$ & 12 & & $17.0$ \\
\phantom{13} 20 & 337 & $430\;25.50$ & & & & \\
& 280 & & $26\;26.97$ & 20 & & $17.5$ \\
& & & $13.25$ & & & $22.0$ \\
\phantom{13} 22 & 337 & $\phantom{0}14\;\phantom{0}3.75$ & $\phantom{0}7\;55.97$ & 12 & $8.1$ & \\
& 281 & & $18\;40.57$ & & $8.1$ & \\
& 449 & $19\;48.53.25$ & & & & \\
& & & $49.25$ & & $7.5$ & \\
\phantom{13} 26 & 394 & $\phantom{0}13\;\phantom{0}\phantom{0}\phantom{0}\phantom{0}$ & $31\;\phantom{0}1.55$ & & & \\
& 620 & & & & & \\
& & $53\;39.25$ & $23\;47.90$ & & & $16.0$ \\
& 452 & & $73$ & & & \\
& 226 & $46$ & $21\;42.54$ & & & $19.0$ \\
& & $51\;31.50$ & & & & \\
& & $\phantom{0}\phantom{0}\phantom{0}\phantom{0}\phantom{0}\phantom{0}23.75$ & $228$ & & & \\
& & & $40\;\phantom{0}2.88$ & & & \\
& & & & & $7.6$ & $19.0$ \\
\bottomrule
\end{tabular}
\end{center}

\newpage

Hieraus finde ich ferner

\begin{center}
\small
\begin{tabular}{c|c|c|c|c|c}
\toprule
& & & & Red.~auf d. & Zenithdistanz \\
& & Refrac- & Parall- & Solstitium & im Solstitium \\
& & tion & axe & & \\
\midrule
Juni 13 & & $1$ & & $-\,0''11$ & \\
\phantom{13} 16 & & $30.73$ & $3.99$ & & $14'\;30''92$ \quad $28^{\circ}\;3'\;48.63$ \\
& & & & & \phantom{0}\phantom{0}\phantom{0}\phantom{0}\phantom{0}\phantom{0}\phantom{0}\phantom{0}\phantom{0}\phantom{0}\phantom{0}\phantom{0}\phantom{0}\phantom{0}\phantom{0}$52''26$ \\
\phantom{13} 18 & & $29''22$ & $3.99$ & $+\,0.34$ & \\
& & $30.33$ & & & $\phantom{0}5\;59.44$ \\
\phantom{13} & & $29.67$ & $3.99$ & $+\,0.26$ & $21\;10.43$ \\
& & & & & \phantom{0}\phantom{0}\phantom{0}\phantom{0}\phantom{0}\phantom{0}\phantom{0}\phantom{0}\phantom{0}\phantom{0}\phantom{0}\phantom{0}\phantom{0}\phantom{0}\phantom{0}$21.96$ \\
\phantom{13} 20 & & $29.25$ & $4.01$ & $+\,0.39$ & \\
\phantom{13} & & $28.99$ & $3.99$ & $+\,0.38$ & $03\;58.31$ \\
& & & & & \phantom{0}\phantom{0}\phantom{0}\phantom{0}\phantom{0}\phantom{0}\phantom{0}\phantom{0}\phantom{0}\phantom{0}\phantom{0}\phantom{0}\phantom{0}\phantom{0}\phantom{0}$23.71$ \\
\phantom{13} 24 & & $29.17$ & $3.99$ & $+\,0.37$ & $\phantom{0}0\;\phantom{0}4.75$ \\
\phantom{13} & & & & & \phantom{0}\phantom{0}\phantom{0}\phantom{0}\phantom{0}\phantom{0}\phantom{0}\phantom{0}\phantom{0}\phantom{0}\phantom{0}\phantom{0}\phantom{0}\phantom{0}\phantom{0}$55.30$ \\
\phantom{13} 21 & & $28.91$ & $3.99$ & $+\,0.30$ & $\phantom{0}0\;\phantom{0}1.81$ \quad $53.46$ \\
\phantom{13} & & & & & \phantom{0}\phantom{0}\phantom{0}\phantom{0}\phantom{0}\phantom{0}\phantom{0}\phantom{0}\phantom{0}\phantom{0}\phantom{0}\phantom{0}\phantom{0}\phantom{0}\phantom{0}$51.41$ \\
\phantom{13} & & & & & \phantom{0}\phantom{0}\phantom{0}\phantom{0}\phantom{0}\phantom{0}\phantom{0}\phantom{0}\phantom{0}\phantom{0}\phantom{0}\phantom{0}\phantom{0}\phantom{0}\phantom{0}$51.77$ \\
\phantom{13} 22 & & $28.89$ & & $-\,0.06$ & $\phantom{0}0\;32.50$ \quad $48.91$ \\
\phantom{13} & & & $3.99$ & $-\,0.08$ & \\
\phantom{13} 29 & & & & & $1\;24.96$ \\
\phantom{13} & & $29.11$ & & & \phantom{0}\phantom{0}\phantom{0}\phantom{0}\phantom{0}\phantom{0}\phantom{0}\phantom{0}\phantom{0}\phantom{0}\phantom{0}\phantom{0}\phantom{0}\phantom{0}\phantom{0}$54.66$ \\
\phantom{13} & & $29.35$ & $3.99$ & $-\,0.05$ & \\
\phantom{13} & & & & & $26\;42.22$ \\
& & & & & \phantom{0}\phantom{0}\phantom{0}\phantom{0}\phantom{0}\phantom{0}\phantom{0}\phantom{0}\phantom{0}\phantom{0}\phantom{0}\phantom{0}\phantom{0}\phantom{0}\phantom{0}$30.75$ \\
\phantom{13} 26 & & $29.47$ & $3.98$ & & $4\;24.19$ \quad $52.10$ \\
\phantom{13} & & $29.12$ & $3.98$ & $-\,0.88$ & \\
& & & & & \phantom{0}\phantom{0}\phantom{0}\phantom{0}\phantom{0}\phantom{0}\phantom{0}\phantom{0}\phantom{0}\phantom{0}\phantom{0}\phantom{0}\phantom{0}\phantom{0}\phantom{0}$53.49$ \\
& & & & & \phantom{0}\phantom{0}\phantom{0}\phantom{0}\phantom{0}\phantom{0}\phantom{0}\phantom{0}\phantom{0}\phantom{0}\phantom{0}\phantom{0}\phantom{0}\phantom{0}\phantom{0}$51.62$ \\
\phantom{13} & & $30.05$ & $4.01$ & & $11\;57.68$ \quad $49.96$ \\
& & & & $15\;17.82$ & \phantom{0}\phantom{0}\phantom{0}\phantom{0}\phantom{0}\phantom{0}\phantom{0}\phantom{0}\phantom{0}\phantom{0}\phantom{0}\phantom{0}\phantom{0}\phantom{0}\phantom{0}$51.61$ \\
& & $29.60$ & $4.01$ & & \\
\bottomrule
\end{tabular}
\end{center}

\begin{center}
Im Mittel aus 214 Beobachtungen \quad $28^{\circ}\;3'\;52''15$ \quad $51''86$
\end{center}

Meine Wahrscheinlichkeitstheorie gibt mir den wahrscheinlichen Fehler dieses Resultats $= \pm\,0''32$.

Nach den Mail"ander Ephemeriden (woraus auch Hr.~Capitain \textsc{Schumacher} bei der Reduction seiner oben
erw"ahnten Beobachtungen die Sonnendeclinationen entlehnte) ist die scheinbare Schiefe der Ekliptik
$= 23^{\circ}\;27'\;52''9$, also Polh"ohe $51^{\circ}\;31'\;45''05$ oder auf den Mittelpunkt reducirt
\[51^{\circ}\;31'\;45''21 \quad [+0.472 \cdot Q \text{ Handschriftliche Aufzeichnung.}]\]
um $4''09$ kleiner als aus dem Polarstern. Nach \textsc{Bessel}'s Schiefe der Ekliptik w"are der Unterschied
noch anderthalb Secunden gr"osser.

Ich muss noch bemerken, dass der Thermometerstand am Barometer nicht mit aufgezeichnet ist.
Bei so m"assigen Zenithdistanzen kann die Vernachl"assigung der Reduction auf eine Normaltemperatur
oder auf die Temperatur der Luft im Freien nur einen sehr kleinen Bruch von einer Secunde Unterschied
geben. Sollte jemand die Beobachtungen nach andern Refractionstafeln und mit R"ucksicht auf eine solche
Refraction berechnen wollen, so wird es wenig fehlen, wenn er jedesmal bei den Sonnenbeobachtungen
die Temperatur des Quecksilbers $2^{\circ}$ niedriger, bei den Polarsternbeobachtungen $2^{\circ}$ h"oher voraussetzt
als die Temperatur im Freien. Zu dieser Bemerkung veranlasst mich die an sich sehr gegr"undete Erinnerung
des Hrn.~Prof.~\textsc{Pasquich} in dem mir eben jetzt eingeh"andigten neuen Heft der Zeitschrift.

\hspace{\fill}[Handschriftlich: $51^{\circ}\;30'\;48''26 - 0.966 \cdot Q$ Wintersolst.~1817.]

\newpage

\subsection*{GAUSS AN BODE.}
\markboth{GAUSS AN BODE.}{}

\begin{center}
\rule{0.6\textwidth}{0.4pt}\\[0.5em]
\textit{Astronomisches Jahrbuch f"ur 1820. Seite 201., 205. Berlin 1817.}\\[0.3em]
\rule{0.6\textwidth}{0.4pt}
\end{center}

\begin{flushright}
G"ottingen, den 17.~M"arz 1817.
\end{flushright}

Da die Theorie der \textsc{Pallas}bewegung hinl"anglich ausgebildet ist, so habe ich Hrn.~Dr.~\textsc{Tittel} aus
Erlau, welcher sich anderthalb Jahre hier aufgehalten, veranlasst, die Ephemeride f"ur 1818..1819 schon
jetzt zu berechnen. Ich habe das Vergn"ugen, Ihnen dieselbe hier f"ur Ihr n"achstes Jahrbuch in Abschrift
beizulegen.

Ist der Himmel im vorigen Winter bei Ihnen den Beobachtungen eben so ung"unstig gewesen,
wie hier? H"ochst selten war die Sonne oder ein Stern zu erblicken. Am ganzen 19.~Nov.~v.~J. schneite
es, und von der Sonnenfinsterniss war nichts zu sehen. Einige Abnahme der Tagshelle war um die
Mitte der Finsterniss zu sp"uren, die jedoch nichts ausserordentliches hatte; wenigstens da es seitdem,
selbst um Mittag, bei Schneegest"ober, auch ohne Sonnenfinsterniss, "ofters eben so dunkel gewesen.
Bei der Mondfinsterniss vom 4.~Dec.~war der Himmel etwas weniger ung"unstig. Anfang und Ende blieben
zwar unsichtbar, allein folgende Phasen habe ich zum Theil zwischen fliegenden Wolken ganz gut beobachtet.

\begin{center}
\textbf{1816. December 4 auf der neuen Sternwarte.}
\end{center}

\begin{center}
\begin{tabular}{l@{\quad\dotfill\quad}r@{\ }r@{\ }r@{\ }r@{\ }r}
\multicolumn{6}{l}{Eintritte:} \\
Tycho 1.~Rand & $8^{\mathrm{h}}$ & $17^{\mathrm{m}}$ & $2^{\mathrm{s}}$ & W.~Z. \\
Tycho 2.~Rand & & $18$ & $44$ & --- \\
Kepler Mitte & & $29$ & $24$ & --- \\
Copernicus 1.~Rand & & $40$ & $28$ & --- \\
Copernicus 2.~Rand & & $45$ & $11$ & --- \\
Aristarch 1.~Rand & & $50$ & $28$ & --- \\
Dionysius Mitte & $9$ & $\phantom{0}0$ & $46$ & --- \\
Censorinus Mitte & & $\phantom{0}5$ & $44$ & --- \\
\end{tabular}
\end{center}

\vspace{1em}

\begin{flushright}
G"ottingen, den 15.~August 1817.
\end{flushright}

Der Bau der hiesigen neuen Sternwarte ist zwar noch nicht ganz, aber doch so weit vollendet,
dass in zwei Zimmern beobachtet werden kann. Bis zur Ankunft und Aufstellung der neuen festen
Meridianinstrumente werden die Beobachtungen mit dem 12~z.~\textsc{Reichenbach}schen Wiederholungskreise
das Erheblichste sein. Ich theile Ihnen davon die Resultate der Beobachtungen des Nordsterns, und der
Sonne im letzten Solstitium mit. Zuv"orderst muss ich bemerken, dass nach meinen geod"atischen Messungen
die Mitte der neuen Sternwarte $6''1$ s"udlicher, und $1^{\mathrm{s}}90$ in Zeit "ostlicher liegt, als die Mitte der alten
Sternwarte. Ferner, dass ich an dem Instrument selbst eine kleine Ab"anderung vorgenommen habe.
Bei meinen fr"ubern Beobachtungen bediente ich mich durchgehends des messingenen Ringes, der an der
Objectivseite des Fernrohrs vorgesteckt wird, und nicht blos das Uebergewicht des prismatischen Oculars,
sondern zugleich die Klammer balancirt. Man hatte vermuthet, dass durch diesen Ring eine Biegung
des Fernrohrs (oder genauer zu reden, eine gr"ossere Biegung der Objectivseite desselben, als durch das
Ocular an der andern Seite erfolgt) hervorgebracht werden k"onnte. Die Folge davon w"urde sein, dass
das Instrument alle Zenithdistanzen zu klein geben m"usste, und man muthmaasste, dass dies vielleicht die
Ursache von dem Unterschiede sei, den man bei Bestimmung der Polh"ohe findet, je nachdem man
Circumpolarsterne oder die Sonne beobachtet. W"are jene Vermuthung gegr"undet, so m"usste der
Unterschied verschwinden, wenn man jenes Gegengewicht wegnimmt, oder vielmehr der Unterschied
m"usste dadurch ins Entgegengesetzte "ubergehen, da entschieden jetzt die Ocularseite des Fernrohrs ein
Uebergewicht hat. Allein meine neuen Beobachtungen, wo jenes Gewicht weggenommen, und dagegen ein
Gewicht an den Alhidadenkreis selbst geschraubt ist, sprechen ganz dagegen, und scheinen mir vollkommen
zu beweisen, dass die Biegung des Fernrohrs durchaus unmerklich ist. Theils geben die neuen
Beobachtungen des Nordsterns und der Sonne wieder ganz denselben Unterschied der Polh"ohe, und in
demselben Sinn, wie die andern; theils ist die absolute Bestimmung der Polh"ohe der neuen Sternwarte
ganz einerlei mit der von der alten Sternwarte "ubertragenen.

\newpage

\section*{NORDSTERN.}
\markboth{NORDSTERN.}{}
\addcontentsline{toc}{section}{Nordstern}

\begin{center}
\textbf{Beobachtungen des Nordsterns in der untern Culmination.}
\end{center}

\begin{center}
\begin{tabular}{c|c|c|c|c}
\toprule
1817 & Anzahl & Wahre Zenithdistanz & & Polh"ohe \\
& & im Meridian & & \\
\midrule
M"arz 28 & 16 & $40^{\circ}\;8'\;13''18$ & & $51^{\circ}\;31'\;49''33$ \\
& & \phantom{0}\phantom{0}\phantom{0}\phantom{0}\phantom{0}13.87 & & \phantom{0}\phantom{0}\phantom{0}\phantom{0}\phantom{0}49.22 \\
April \phantom{0}2 & 22 & & & \phantom{0}\phantom{0}\phantom{0}\phantom{0}\phantom{0}50.13 \\
\phantom{0}\phantom{0}8 & & \phantom{0}\phantom{0}\phantom{0}\phantom{0}\phantom{0}15.86 & & \phantom{0}\phantom{0}\phantom{0}\phantom{0}\phantom{0}47.84 \\
& & \phantom{0}\phantom{0}\phantom{0}\phantom{0}\phantom{0}15.76 & & \phantom{0}\phantom{0}\phantom{0}\phantom{0}\phantom{0}50.03 \\
11 & 26 & \phantom{0}\phantom{0}\phantom{0}\phantom{0}\phantom{0}16.16 & & \phantom{0}\phantom{0}\phantom{0}\phantom{0}\phantom{0}46.75 \\
& & \phantom{0}\phantom{0}\phantom{0}\phantom{0}\phantom{0}20.51 & $51^{\circ}\;31'$ & \\
\phantom{0}\phantom{0} & 12 & & & \\
Mai \phantom{0}7 & & \phantom{0}\phantom{0}\phantom{0}\phantom{0}\phantom{0}23.90 & & \phantom{0}\phantom{0}\phantom{0}\phantom{0}\phantom{0}50.51 \\
\midrule
Mittel & 158 & & & $51^{\circ}\;31'\;49''.60$ \\
\bottomrule
\end{tabular}
\end{center}

Bei Berechnung der Refraction ist \textsc{Carlini}'s Tafel, f"ur die Declination des Nordsterns die \textsc{Bessel}sche
zum Grunde gelegt. Der Beobachtungsplatz liegt $0''12$ n"ordlich von der Mitte; die Polh"ohe von dieser
ist also, aus obigen Beobachtungen $51^{\circ}\;31'\;49''3$.

Meine Beobachtungen des Nordsterns in der untern Culmination von 1813 geben mit denselben
Elementen reducirt die Polh"ohe der alten Sternwarte $51^{\circ}\;31'\;55''48$; die Uebertragung gibt also
dasselbe was die unmittelbaren Beobachtungen geben.

Die Beobachtungen des Sommer-Sonnenstillstandes wurden in diesem Jahre von der Witterung
mehr als je beg"unstigt: der Beobachtungsplatz liegt $0''16$ s"udlich von der Mitte.

\begin{center}
\textbf{1817 Beobachtungen des Sonnenstillstandes.}
\end{center}

\begin{center}
\small
\begin{tabular}{c|c|c|c|c}
\toprule
1817 & Anzahl & Wahre Zenithdistanz & Im Solstitium & \\
& & im Meridian & & \\
\midrule
Juni 13 & & & & \\
\phantom{13} 16 & & & $28^{\circ}\;3'\;52''26$ & \\
& & & \phantom{0}\phantom{0}\phantom{0}\phantom{0}\phantom{0}48.63 & \\
\phantom{13} 18 & 10 & $28^{\circ}\;18'\;\phantom{0}\phantom{0}\phantom{0}\phantom{0}\phantom{0}\phantom{0}\phantom{0}\phantom{0}$ & & \\
\phantom{13} & 16 & \phantom{0}\phantom{0}\phantom{0}\phantom{0}\phantom{0}$6\;19.03$ & & \\
\phantom{13} 20 & 12 & \phantom{0}\phantom{0}\phantom{0}\phantom{0}\phantom{0}$9\;51.80$ & & \\
\phantom{13} & & \phantom{0}\phantom{0}\phantom{0}\phantom{0}\phantom{0}$5\;\phantom{0}8.69$ & & \\
\phantom{13} & 14 & \phantom{0}\phantom{0}$47\;57.18$ & & \phantom{0}\phantom{0}\phantom{0}\phantom{0}\phantom{0}$55.30$ \\
\phantom{13} & 20 & & & \phantom{0}\phantom{0}\phantom{0}\phantom{0}\phantom{0}51.77 \\
\phantom{13} 22 & & & & \phantom{0}\phantom{0}\phantom{0}\phantom{0}\phantom{0}53.46 \\
\phantom{13} & 27 & $29\;\phantom{0}\phantom{0}\phantom{0}\phantom{0}\phantom{0}\phantom{0}\phantom{0}\phantom{0}\phantom{0}$ & & \phantom{0}\phantom{0}\phantom{0}\phantom{0}\phantom{0}48.91 \\
\phantom{13} 21 & & & & \\
\phantom{13} & 10 & \phantom{0}\phantom{0}$43\;54.41$ & & \\
\phantom{13} & 12 & \phantom{0}\phantom{0}\phantom{0}$6\;36.39$ & & \phantom{0}\phantom{0}\phantom{0}\phantom{0}\phantom{0}54.66 \\
\phantom{13} 26 & 16 & \phantom{0}$108\;28.75$ & & \phantom{0}\phantom{0}\phantom{0}\phantom{0}\phantom{0}52.10 \\
\phantom{13} & & \phantom{0}\phantom{0}\phantom{0}20.16 & & \phantom{0}\phantom{0}\phantom{0}\phantom{0}\phantom{0}51.62 \\
\phantom{13} & & & & \phantom{0}\phantom{0}\phantom{0}\phantom{0}\phantom{0}53.49 \\
\phantom{13} & 22 & \phantom{0}\phantom{0}$15\;19\;54.34$ & & \phantom{0}\phantom{0}\phantom{0}\phantom{0}\phantom{0}49.96 \\
\phantom{13} & & \phantom{0}\phantom{0}\phantom{0}16.59 & & \phantom{0}\phantom{0}\phantom{0}\phantom{0}\phantom{0}51.41 \\
\phantom{13} & \phantom{0}8 & & & \phantom{0}\phantom{0}\phantom{0}\phantom{0}\phantom{0}51.61 \\
\midrule
Mittel aus 214 Beob. & & & $28^{\circ}\;3'\;52''15$ & $51''86$ \\
\bottomrule
\end{tabular}
\end{center}

\newpage

\section*{COMET.}
\markboth{COMET.}{}

\begin{center}
\rule{0.6\textwidth}{0.4pt}\\[0.5em]
\textit{V.~\textsc{Lindenau} u.~\textsc{Bohnenberger} Zeitschrift f"ur Astronomie u.~verwandte Wiss.}\\
\textit{Band V. Seite 276\,..\,277. M"arz und April 1818.}\\[0.3em]
\rule{0.6\textwidth}{0.4pt}
\end{center}

Aus den Beobachtungen des Cometen von Dr.~\textsc{Olbers} vom 3.~bis 28.~M"arz habe ich dieser Tage
folgende parabolische Elemente abgeleitet, die von den von \textsc{Olbers} selbst aus seinen Beobachtungen
vom 3.~und 13.~M"arz und der mir noch unbekannten Marseiller vom 18.~Januar erhaltenen, nur wenig
abweichen.

\begin{center}
\begin{tabular}{l@{\dotfill}l}
Durchgang durch das Perihel 1818.~Februar 26,8991 & G"ottinger Zeit. \\
Logarithm.~des Abstandes im Perihel & $0.07937$ \\
Neigung der Bahn & $90^{\circ}\;0'\;0''$ \\
Aufsteigender Knoten & $7\;\phantom{0}5\;\phantom{0}\phantom{0}\phantom{0}$ \\
L"ange des Perihels, die Bewegung als rechtl"aufig betrachtet & $183\;22\;58$ \\
\end{tabular}
\end{center}

Die senkrechte Lage der Bahn gegen die Ekliptik ergab sich ungezwungen aus den Beobachtungen,
wird aber wohl nach fortgesetzten Beobachtungen einige Modification erhalten. Dann erst wird man
sehen, ob die Bewegung in Beziehung auf die Ekliptik rechtl"aufig oder r"uckl"aufig ist. Ich habe zur
Erleichterung der Beobachtungen aus diesen Elementen folgende Ephemeride berechnet. Die Zeit ist
$14^{\mathrm{h}}\;20^{\mathrm{m}}$ in G"ottingen. Die Columne f"ur die Lichtst"arke setzt erborgtes Licht voraus, $=1$ in der
Distanz 1 von Erde und Sonne.

\begin{center}
\small
\begin{tabular}{c|c|c|c}
\toprule
1818 & AR. & Decl. & Lichtst"arke \\
\midrule
April 10 & $297^{\circ}\;50'$ & $12^{\circ}\;57'$~N. & $0.503$ \\
\phantom{10} 14 & $296\;27$ & $\phantom{0}6\;30$ & $0.471$ \\
\phantom{10} 18 & $298^{\circ}\;47'58''$ & $\phantom{0}0\;43$ & \\
\phantom{10} & $294\;\phantom{0}\phantom{0}\phantom{0}\phantom{0}\phantom{0}$ & & \\
\phantom{10} 22 & $292\;44$ & $\phantom{0}6\;41$ & \\
\phantom{10} 26 & $290\;17$ & $\phantom{0}2\;46$ & $0.596$ \\
& $283\;53$ & & $0.568$ \\
Mai \phantom{0}4 & $287\;21$ & $3\;48$~S. & $0.616$ \\
\phantom{0}\phantom{0} \phantom{0}8 & $279\;47$ & $10\;58$ & \\
\phantom{0}\phantom{0} 12 & $275\;\phantom{0}4$ & $15\;33$ & $0.621$ \\
\bottomrule
\end{tabular}
\end{center}

L"anger werden wegen des Mondscheins die Beobachtungen wohl nicht fortgesetzt werden k"onnen,
und wenn dieser aufh"ort, wird der Comet f"ur uns wenigstens zu weit nach S"uden fortger"uckt sein.

Ich selbst habe den Cometen bis jetzt erst einmal beobachtet, und mit einem Stern 9ter Gr"osse
verglichen, der in der Hist.~C\kern-0.5pt{}el. nicht vorkommt. Ich behalte mir vor, dessen Ort k"unftig zu bestimmen;
ungef"ahr wird er $300^{\circ}\;3'$~AR., $15^{\circ}\;9'$ Declination haben.

\begin{center}
\begin{tabular}{l@{\quad}l}
April 1. & $15^{\mathrm{h}}\;45^{\mathrm{m}}\;28^{\mathrm{s}}$ Sternzeit, der Comet folgt dem Stern $14^{\mathrm{s}}98$ in Zeit \\
& $15\;44\;56$ \phantom{Sternzeit,} der Comet ist $4'\;22''$ s"udlicher als der Stern.
\end{tabular}
\end{center}

Die Declination m"ochte wohl betr"achtlich weniger genau sein, als die Rectascension.

\hspace{\fill}[Handschrift: Beobachtungen dieses Sterns \quad 1819 Aug.~27 \quad $20^{\mathrm{h}}\;0^{\mathrm{m}}\;11^{\mathrm{s}}49$\\
\null\hfill Sept.~\phantom{0}4 \quad \phantom{0}\phantom{0}\phantom{0}\phantom{0}\phantom{0}$11.48$\\
\null\hfill \phantom{0}\phantom{0}\phantom{0}\phantom{0}\phantom{0}\phantom{0}7 \quad \phantom{0}\phantom{0}\phantom{0}\phantom{0}\phantom{0}$11.38$\\
\null\hfill \phantom{0}\phantom{0}\phantom{0}\phantom{0}\phantom{0}24 \quad \phantom{0}\phantom{0}\phantom{0}\phantom{0}\phantom{0}$11.30$]

\vspace{1em}

\begin{center}
\rule{0.6\textwidth}{0.4pt}\\[0.5em]
\textit{G"ottingische gelehrte Anzeigen. St"uck 60. Seite 593\,..\,594. 1818.~April 13.}\\[0.3em]
\rule{0.6\textwidth}{0.4pt}
\end{center}

Zur Erleichterung der Beobachtung des vom Hrn.~\textsc{Pons} Ende vorigen Jahrs entdeckten Cometen
hat Hr.~Hofr.~\textsc{Gauss} f"ur die noch "ubrige Zeit seiner Sichtbarkeit den Lauf desselben im Voraus berechnet;
wir glauben durch die Bekanntmachung dieser Resultate den Beobachtern einen Dienst zu erweisen.
Die parabolischen Elemente, nach denen diese Ephemeride berechnet ist, gr"unden sich auf die Beobachtungen
des Hrn.~Doctor \textsc{Olbers} vom 3.~bis 28.~M"arz und sind folgende:

\begin{center}
\begin{tabular}{l@{\quad}l@{\quad}l}
Zeit des Durchgangs durch die Sonnenn"ahe & Aufsteigender Knoten & $7^{\circ}\;0'$ \\
1818 Febr.~26.\ $21^{\mathrm{h}}\;34^{\mathrm{m}}\;42^{\mathrm{s}}$ Zeit im Mer.~von G"ottingen & Neigung der Bahn & $90$ \\
Kleinster Abstand von der Sonne \dotfill $1.20053$ & L"ange der Sonnenn"ahe & $183$ \\
\end{tabular}
\end{center}

Die k"unftige Verbesserung dieser Elemente nach mehrern Beobachtungen wird erst lehren ob die
Neigung merklich vom rechten Winkel abweicht, und in welchem Sinn; bis jetzt bleibt unbestimmt,
ob die Bewegung in Beziehung auf die Ekliptik rechtl"aufig oder r"uckl"aufig ist. Die Zeit der Ephemeride
ist $14^{\mathrm{h}}\;20^{\mathrm{m}}$ im Meridian von G"ottingen. [Die Ephemeride ist in dem vorhergehenden Aufsatze abgedruckt.]

Obgleich der Comet schon weit "uber seine Sonnenn"ahe hinaus ist, so wird doch sein Licht wegen der
noch abnehmenden Entfernung von der Erde noch immer zunehmen, so dass es am 12.~Mai doppelt so gross
sein wird als am 13.~M"arz; allein da dasselbe an sich so ausserordentlich schwach ist, so wird der Comet
doch nur durch gute Fernr"ohre zu beobachten sein.

Die auf der hiesigen Sternwarte angestellten Beobachtungen dieses Cometen werden bei einer andern
Gelegenheit bekannt gemacht werden.

\newpage

\section*{BORDA'S KREIS.}
\markboth{BORDA'S KREIS.}{}
\addcontentsline{toc}{section}{Borda's Kreis}

\begin{center}
\rule{0.6\textwidth}{0.4pt}\\[0.5em]
\textit{V.~\textsc{Lindenau} u.~\textsc{Bohnenberger}, Zeitschrift f"ur Astronomie u.~verwandte Wiss.}\\
\textit{Band V. Seite 198\,..\,211. M"arz und April 1818.}\\[0.3em]
\rule{0.6\textwidth}{0.4pt}
\end{center}

Ihrem Wunsche zufolge, verehrtester Herr Professor, habe ich das Vergn"ugen, Ihnen hier diejenige
Methode mitzutheilen, deren ich mich zur Parallelstellung der Gesichtslinie mit der Kreisebene bedient
habe, und die, wie eine oft wiederholte Erfahrung mich gelehrt hat, einer sehr grossen Sch"arfe f"ahig ist.

Um mich kurz und bestimmt ausdr"ucken zu k"onnen, bezeichne ich die verschiedenen Axen des Instruments
auf folgende Weise. Diejenige Axe, die bei H"ohenmessungen vertikal gestellt werden muss,
und die man auch bei jedem andern Gebrauch wenigstens nicht viel von dieser Lage abweichen l"asst,
nenne ich die \textit{erste Axe}. Die Bewegung um diese Axe wird auf dem Azimuthkreise gemessen; da jedoch
derer nur Einen Index hat, dessen Vernier nur Minuten angibt, obwohl noch Theile von Minuten sich
sch"atzen lassen, und man nicht sicher ist, ob nicht eine kleine Excentricit"at vorhanden ist, so kann eine
Berichtigung der Gesichtslinie, vermittelst dieses Azimuthalkreises (indem man bei zwei entgegengesetzten
vertikalen Lagen des Hauptkreises das Azimuth eines irdischen Gegenstandes, n"othigenfalls mit R"ucksicht
auf seine Parallaxe, wenn er nicht sehr entfernt ist, bestimmt) noch nicht alle erreichbare oder w"unschenswerthe
Genauigkeit geben; inzwischen mag man sich dieses bequemen und einfachen Verfahrens zur ersten
vorl"aufigen Berichtigung bedienen.

Nahe senkrecht zur ersten Axe ist diejenige, die ich die \textit{zweite} nenne; diese ist also jederzeit nahe
horizontal. Durch Drehung um dieselbe wird der Kreis in die vertikale, schiefe oder horizontale Lage
gebracht. An den \textsc{Reichenbach}schen Kreisen ist keine Theilung, um die Gr"osse dieser Drehung zu
messen. F"ur die meisten Beobachtungen ist dies auch "uberfl"ussig; es gibt aber doch einige F"alle, wo es
w"unschenswerth w"are, namentlich w"urde man ohne M"uhe Winkel zwischen Fixsternen und irdischen
Objecten bei Tage messen k"onnen, was jetzt entweder gar nicht oder nur mit grosser Schwierigkeit zu
bewerkstelligen ist.

Wiederum fast senkrecht zur zweiten Axe ist die, um welche sich der ganze Kreis dreht, welche
Bewegung am Trommelkreise gemessen wird. Da mit letzterer Axe, im mathematischen Sinn, diejenige,
um welche der Alhidadenkreis sich bewegt, zusammenf"allt, so werden diese beiden Axen nur als Eine
betrachtet, die ich die \textit{dritte} nenne. Diese dritte Axe ist auf der Kreisebene genau senkrecht, und wenn
$\alpha$ den Winkel der Gesichtslinie mit der Kreisebene bedeutet (welchen ich positiv setze, wenn die Objectivseite
weiter von der Kreisfl"ache abliegt, als die Okularseite), so ist $90^{\circ} - \alpha$ der Winkel der Gesichtslinie mit der
dritten Axe, jene nach der Objectivseite, diese abw"arts von der getheilten Fl"ache genommen.

F"ur den Gebrauch des Instruments ist nicht nothwendig, dass die zweite Axe mit der Kreisebene
genau parallel, und zur ersten Axe genau senkrecht sei; die Unterschiede werden jedoch gewiss klein sein,
und eine oft wiederholte Erfahrung an dem hiesigen 12z"olligen \textsc{Reichenbach}schen Kreise hat mich gelehrt,
dass sie Jahre lang, und selbst, nachdem das Instrument auseinandergenommen gewesen, unver"anderlich
bleiben. Ich nenne $b$ die Neigung der zweiten Axe gegen die Kreisebene, positiv genommen, wenn der
linke Arm\footnote{insofern man hinter dem senkrecht vorausgesetzten Kreise sich befindet.} von jeener dem
Kreise n"aher ist. Endlich nenne ich $c$ die Neigung der zweiten Axe gegen eine auf der ersten Axe
senkrechten Ebene (wof"ur man die Ebene des Azimuthalkreises halten kann), positiv genommen, wenn der
linke Arm mehr "uber dieselbe erhaben ist.

Es seien noch $M$ und $N$ die Schnitte der dritten Axe mit der zweiten und mit der Gesichtslinie
und $O$ ein gut zu sehendes Object. Die Entfernung $MN$ hat an jedem Kreise eine bestimmte, durch
Abmessung oder auf andere Weise anzugebende Gr"osse, an dem hiesigen $\frac{1}{2}$~Meter. Insofern die Gesichtslinie
auf das Object $O$ gerichtet ist, bezeichne ich den Winkel $MON$ mit $p$; die genaue Formel daf"ur ist
\[\sin p = \frac{MN}{MO \cdot \cos\alpha}\]
wof"ur man jederzeit $p = \frac{MN}{MO} \cdot 206265''$ annehmen darf. Es wird nicht schwer sein, $MO$ so genau zu
erfahren, dass $p$ auf die Sekunde genau bekannt wird. Bei Objecten in bedeutender Entfernung wird
$p$ ganz unmerklich; dies ist aber kein Grund, sie n"ahern bei der folgenden Berichtigungsmethode vorzuziehen,
wenn letztere bessere Zielpuncte abgibt.

Ich setze voraus, dass bei allen hier zu beschreibenden Operationen der Kreis selbst um die dritte
Axe gar nicht bewegt wird, also die Hemmung an der Trommel fest angezogen ist, und die Steilschraube
an derselben gar nicht ber"uhrt wird. Man kann, wenn man will, auch die Pressmutter am hintern Theile
der dritten Axe etwas st"arker als gew"ohnlich anziehen.

Ich nenne K"urze halber $I$ den Winkel zwischen einer Ebene durch die zweite und dritte Axe mit der
Ebene durch die erste und zweite, welcher mithin die Gr"osse der Drehung des Kreises um die zweite Axe
bestimmt. Der Bau des Instruments verstattet nur Werthe dieses Winkels von 0 bis $90^{\circ}$, und etwas auf
beiden Seiten "uber diese Grenzen hinaus; dass an den \textsc{Reichenbach}schen Kreisen keine Theilung zur Messung
dieses Winkels sich befindet, habe ich schon oben bemerkt.

Das Wesentliche meiner Methode beruht darauf, dass, wenn bei irgend einer Lage des Instruments
die Gesichtslinie auf ein Object gerichtet ist, durch Verbindung einer Drehung um die zweite Axe mit
einer Drehung des Alhidadenkreises um die dritte, die Gesichtslinie zum zweitenmale auf das Object
gef"uhrt werden kann. Ich nehme an, es sei $I = \iota$, wenn die dritte Axe mit der zweiten und der geraden
Linie $MO$ in Eine Ebene gebracht ist, der Index des Alhidadenkreises hingegen zeige $k$, wenn die
Gesichtslinie mit der zweiten und der dritten Axe in Einer Ebene sich befindet, so dass die Gesichtslinie
mit der zweiten Axe (linker Arm) auf der Ebene des Kreises einerlei Projection haben.

\newpage

Man "uberzeugt sich leicht, dass wenn bei der Richtung der Gesichtslinie auf das Object der Index des
Alhidadenkreises auf $k - x$ steht, und $I = \iota - \eta$ ist, die Gesichtslinie wiederum auf dieses Object gerichtet
sein wird, wenn $I = \iota + \eta$ gemacht, und der Index des Alhidadenkreises auf $k + x$ gestellt wird; diese
zweite Stellung wird nat"urlich nur dann physisch m"oglich sein, wenn der zweite Werth von $I$ innerhalb
der Grenzen f"allt, die die Einrichtung des Instruments zul"asst. Es l"asst sich ferner leicht darthun, dass
nach aller Strenge folgende Gleichung Statt findet:
\[\sin b \cdot \cos x = \cos b \cdot \tan(\alpha + p) + \cot \eta \cdot \sin x\]
und dass der Cosinus des Winkels der zweiten Axe mit der Linie $MO$
\[= \sin b \cdot \sin(\alpha + p) + \cos b \cdot \cos(\alpha + p) \cdot \cos x\]
wird. Hieraus folgert man leicht, dass $\eta$ entweder sehr klein, oder sehr nahe $180^{\circ}$ sein wird; im ersten
Fall wird der gedachte Winkel sehr klein, im zweiten nahe $180^{\circ}$ sein, d.~i.~im ersten Fall wird das Object
nahe in der Richtung des linken Arms der zweiten Axe, im zweiten nahe in der Richtung des rechten liegen m"ussen.

Man hat daher hinreichend genau

\begin{center}
\begin{tabular}{ll@{\quad}l@{\quad}l}
im ersten Fall: & $b - \alpha - p = \sigma - k = k - h = x(\iota - h)$ & $\tan\eta$ \\
im zweiten Fall: & $b + \alpha + p = \sigma' - k = k - h' = x(\iota' - h')$ & $\eta = 180^{\circ} - \tau$
\end{tabular}
\end{center}

Durch zwei solche Versuche erh"alt man daher, insofern man sich eine Kenntniss von $\eta$ verschaffen
kann, die Werthe von $b - \alpha - p$ und von $b + \alpha + p$, also folglich von $b$ und $\alpha + p$, oder, insofern
$p$ als bekannt anzusehen ist, von $\alpha$; zugleich erh"alt man offenbar aus jedem solchen Versuch (ohne
$\eta$ n"othig zu haben) eine Bestimmung von $k$. Da $\eta$ hiebei willk"uhrlich ist, so ist es am vortheilhaftesten,
es so gross wie m"oglich zu machen. Indessen w"urde man doch nicht viel "uber $45^{\circ}$ hinaus gehen k"onnen,
und da man an \textsc{Reichenbach}'schen Kreisen die Werthe von $I$ nicht ablesen kann, so muss man sich
auf solche beschr"anken, wof"ur man die Stellung auf anderm Wege erhalten kann. Ich werde unten zeigen,
wie man die Stellung, wof"ur $I = 0$ ist, d.~i.~wo die erste, zweite und dritte Axe in Einer Ebene liegen,
mit gr"osster Sch"arfe effectuiren kann. Man kann ferner die Kreisebene mit der ersten Axe parallel
machen, indem man auf bekannte Weise beide vertical stellt; man kann diese beiden Stellungen durch
zwei auf dem Quadranten, an dem sich die Hemmung f"ur die Drehung um die zweite Axe befindet,
an der Kante der einen der St"utzen, die die zweite Axe tragen, gezogene Striche bezeichnen, um sie
immer schnell und hinreichend genau wieder finden zu k"onnen. Ich werde die erste Stelle K"urze halber
die \textit{horizontale}, die zweite die \textit{vertikale} nennen. F"ur letztere ist nach aller Sch"arfe
\[\cos I = -\tan b \cdot \tan c.\]
Es wird daher $I$ nicht merklich von $90^{\circ}$ verschieden sein, und da der Unterschied $\iota - \eta$ gesetzt ist,
so wird genau
\[x = \sqrt{\frac{-\tan \eta}{1 + \tan b \cdot \tan c}} \cdot \frac{\cos(b - c)}{\cos(b + c)}\]
also nur unmerklich von $\eta$ verschieden sein.

\newpage

Hieraus folgt also, dass wenn bei dem ersten Versuch, wo das Object nahe in der Richtung des
linken Arms der zweiten Axe sich befunden hat, bei der horizontalen Stellung des Kreises der Index
des Alhidadenkreises $h$, bei der vertikalen $u$ gezeigt hat, im zweiten Versuch hingegen, wo das Object
sich nahe in der Richtung des rechten Arms der zweiten Axe befand, der Index des Alhidadenkreises
respectiv $180^{\circ} + h'$, $180^{\circ} + u'$ zeigte, man haben wird:
\begin{align*}
b - \alpha - p &= \sigma - k = k - h = x(\iota - h) \\
b + \alpha + p &= u' - h' = k - h' = x(\iota' - h')
\end{align*}
folglich
\begin{align*}
\alpha &= \tfrac{1}{2}(u' - h' - u + h) - p \\
b &= \tfrac{1}{2}(u' - h' + u - h)
\end{align*}
Zugleich hat man f"ur $k$ die doppelten Werthe
\[\tfrac{1}{2}(u' + h) \quad \text{und} \quad \tfrac{1}{2}(u + h')\]
deren nahe Uebereinstimmung zur Best"atigung der Richtigkeit der Operationen dienen wird.

Der Gang bei meinem Verfahren ist demnach folgender. Das Object, welches ich dabei w"ahle,
muss einen scharfen Zielpunct darbieten und h"ochstens einige Grade vom Horizont entfernt sein. Ich
stelle das Instrument so auf, dass der eine Fuss desselben mit der ersten Axe und dem Object n"ahe in
Einer Ebene liegt, in welche ich sodann auch nach dem Augenmaass durch Azimuthalbewegung die zweite
Axe bringe, sodass zuerst deren linker Arm dem Object zugekehrt ist. Hiern"achst stelle ich den Kreis
horizontal, und bringe durch die Schraube an dem erw"ahnten Fusse und durch Alhidadenbewegung die
Gesichtslinie auf das Object. Sodann wird der Kreis vertical gestellt, und die Gesichtslinie durch
Azimuthal- und Alhidadenbewegung von neuem auf das Object gebracht. Hiedurch ist die erste grobe
Stellung vollendet. Der Kreis wird von neuem horizontal gestellt, und falls die Alhidadenbewegung nicht
hinreicht, um die Gesichtslinie auf das Object zu f"uhren, an der Fussschraube nachgeholfen. Dann wird
der Kreis wiederum vertical gestellt, und die Gesichtslinie durch Alhidadenbewegung, n"othigenfalls mit
Nachh"ulfe an der Azimuthschraube, auf das Object gebracht. Mit diesen beiden Operationen wird
abwechselnd so lange fortgefahren, bis die blosse Alhidadenbewegung in horizontaler Lage zum scharfen
Pointiren hinreicht. Die Ablesungen am Alhidadenkreise geben $h$ und $u$.

Hiern"achst wird nach der letzten Verticalstellung die Azimuthalhemmung gel"ost, der Kreis $180^{\circ}$
um die erste Axe bewegt, so dass jetzt der rechte Arm der zweiten Axe n"ahe auf das Object gerichtet
ist; durch Alhidadenbewegung (die n"ahe $180^{\circ}$ betragen wird) und feine Stellung an der Azimuthschraube
die Gesichtslinie scharf auf das Object gebracht; auch hier lese ich den Index des Alhidadenkreises ab,
welcher $180^{\circ} + h'$ geben mag. Dies ist der Anfang der zweiten Operationsreihe. Der Kreis wird
horizontal gelegt, und durch Stellung an der Fussschraube und Alhidadenbewegung pointirt; bei der
neuen Verticalstellung wird Azimuthal- und Alhidadenbewegung angewandt. So wird wieder wechselsweise
fortgefahren, bis die blosse Alhidadenbewegung zureicht, das Object in beiden Lagen zu treffen.
Die Ablesungen am Alhidadenkreise geben hier $180^{\circ} + h'$ und $180^{\circ} + u'$. Die Messungen sind hiermit
vollendet.

Die vorhin angef"uhrte Ablesung $180^{\circ} + h'$ dient noch dazu, auch die Gr"osse $c$ zu bestimmen. Offenbar
wird, bei verticaler Stellung des Kreises, die Gesichtslinie mit der ersten und dritten Axe in Einer
Ebene, und zwar nach oben zugerichtet, sein, wenn der Index des Alhidadenkreises auf $k + V - 90^{\circ}$
steht, oder die Projection der Gesichtslinie auf der Kreisebene wird dann mit der ersten Axe parallel sein.
Die Projection der zweiten Axe (linker Arm) auf eben dieser Ebene macht also mit dem Radius, der der
ersten Axe parallel ist, den Winkel $h + 90^{\circ} - (u + V)$. Es ist daher
\[\sin c = \cos b \cdot \sin(h + V - u - k)\]
wof"ur man
\[c = \cos b \cdot (h + V - u - k)\]
setzen darf.

Dieses Verfahren f"uhrt also zur Kenntniss des Fehlers $\alpha$ in Sekunden. Will man ihn wegschaffen,
so muss das Fadenkreuz, insofern $\alpha$ positiv ist, abw"arts vom Kreise geschoben werden, oder, wenn
man sich des prismatischen Okulars bedient, in der Richtung vom Objectiv zum Okular. Um die Gr"osse
der Verschiebung anschaulich vor sich zu haben, kann man sich entweder einer Sch"atzung des Verh"altnisses
zu einem schicklichen Object von bekanntem Durchmesser nach dem Augenmaass bedienen, oder zuerst die
Gesichtslinie scharf auf ein feines Object bringen, und dann den Alhidadenkreis um die Gr"osse $\alpha$
verschieben, wo dann das Fadenkreuz soviel verr"uckt werden muss, bis das Object in der Diagonalen
erscheint. Dass man bei dieser delikaten Operation vorsichtig zu Werke gehen und alle Verr"uckung des
Kreises selbst, w"ahrend derselben, vermeiden muss, versteht sich von selbst. Immer bietet die Wiederholung
der vorstehenden Methode das Mittel dar, sich zu versichern, ob man seinen Zweck erreicht habe.

Es bleibt mir jetzt noch "ubrig, zu zeigen, wie die horizontale Stellung des Kreises, wo $I = 0$,
zu effectuiren sei. Insofern $c$ und $b$ nicht betr"achtlich sind, kommt es freilich so genau nicht darauf
an, und man k"onnte sich mit dem Augenmaasse begn"ugen. Allein durch folgendes Verfahren l"asst sich
dieses mit "ausserster Sch"arfe erreichen.

Nachdem bei festem Stande der Trommel $k$ bestimmt ist, lege man zuerst nach dem Augenmaass
den Kreis horizontal, stelle den Index des Alhidadenkreises auf $k + l$, wo $l$ einen beliebigen Winkel,
am besten $90^{\circ}$, bedeutet, und bringe durch Azimuthal- und Fussbewegung ein gut zu sehendes Object
auf die Gesichtslinie. Hierauf stelle man den Index des Alhidadenkreises auf $k - l$; hat man die Horizontalstellung
zuf"allig genau getroffen, so wird durch blosse Azimuthalbewegung die Gesichtslinie wieder
genau auf das Object gef"uhrt werden k"onnen, sonst aber dar"uber oder darunter weggehen. Nachdem
man dann durch die Azimuthalbewegung dies Object der Gesichtslinie so nahe wie m"oglich (d.~i.~auf
den der dritten Axe parallelen, oder bei prismatischem Okular auf den mit dem Fernrohr parallelen
Faden) gef"uhrt hat, bringe man durch Drehung des Kreises um die zweite Axe das Object der Gesichtslinie
um die H"alfte n"aher und corrigire die andere H"alfte durch eine Fussschraube, n"othigenfalls mit einer
Nachh"ulfe der Azimuthschraube, so dass das Object wieder scharf auf der Gesichtslinie steht. Um die
Genauigkeit der Halbirung zu pr"ufen, stelle man von neuem den Index des Alhidadenkreises auf $k + l$
und versuche, ob jetzt die Azimuthalbewegung allein das Object genau auf die Gesichtslinie zur"uckf"uhrt,
wo nicht, so corrigirt man wieder die H"alfte des jetzt gewiss viel kleinern Fehlers an der Schraube
ohne Ende, wodurch die feine Bewegung um die zweite Axe geschieht. Mit dieser Operation f"ahrt man
so lange fort, bis man vollkommen befriedigt ist.

W"are die erste, nach dem Augenmaasse gemachte, Stellung gar zu schlecht gerathen, so w"urde
es schwer sein, die erste Halbirung leidlich genau zu machen, weil die Azimuthalbewegung, insofern $l$
gross angenommen wurde, das Object gar nicht in das Gesichtsfeld bringen w"urde. Um sich also viel
Probiren zu ersparen, nehme man Anfangs $l$ lieber klein, etwa $= 10^{\circ}$, und bringe den Kreis nach
vorgeschriebener Methode in die verlangte Stellung, welche denn durch Wiederholung des Verfahrens
bei einem grossen Werth von $l$ fein berichtigt werden kann.

Wenn an einem Instrument zuf"allig $l = c$ ist, so wird die Kreisebene bei horizontaler Stellung
genau auf die erste Axe senkrecht sein, und ein Object, welches einmal auf der Gesichtslinie ist, wird
bei jeder ver"anderten Alhidadenstellung durch blosse Azimuthalbewegung auf dieselbe zur"uckgebracht.
Sind aber $b$ und $c$ ungleich, so "uberzeugt man sich leicht, dass dies unm"oglich ist; die Kreisebene kann
dann im eigentlichen Sinne nicht genau horizontal gestellt werden, wenn nicht die erste Axe etwas geneigt
wird. Bei dem hiesigen \textsc{Reichenbach}schen Kreise ist $c = +5'\;40''$, $l = +28''$ im Mittel aus mehrern
Versuchen.

Uebrigens ist schon oben bemerkt, dass die Bestimmung von $k$ von der Kenntniss bestimmter
Werthe von $I$ ganz unabh"angig ist. Man braucht bloss, ohne Azimuthal- und Trommelschraube zu
ber"uhren, durch Verbindung von Alhidadenbewegung und Drehung um die zweite Axe ein Object zweimal
auf die Gesichtslinie zu f"uhren und das Mittel aus beiden Stellungen des Index des Alhidadenkreises zu
nehmen. Am zweckm"assigsten ist es hiebei, es so einzurichten, dass die Drehung um die zweite Axe so
gross wie m"oglich wird. Um dies zu erhalten, mache ich vorl"aufig die oben gelehrte grobe Stellung, und
wenn dann, bei ungef"ahr horizontaler Stellung des Kreises, das Object genau auf der Gesichtslinie ist,
drehe ich den Kreis um die zweite Axe, bis die Bogenbewegung, die das Object hierdurch im Gesichtsfelde
erh"alt, dasselbe wieder auf den Faden f"uhrt, der mit der Kreisebene parallel ist (bei prismatischem
Okular auf den, welcher diesen repr"asentirt, d.~i.~senkrecht zum Fernrohr steht), und bringe durch
Alhidadenbewegung dann das Object genau auf das Kreuz. Bliebe bei der Drehung um die zweite Axe
das Object genau auf dem Kreuz, so w"are dies ein Beweis, dass das Object in der zweiten Axe selbst
l"age, dass folglich $b = \alpha + p$, und die Stellung des Index selbst $= k$.

Ich will bei dieser Gelegenheit noch des Verfahrens erw"ahnen, dessen ich mich mit Vortheil bediene,
um die Kreisebene genau vertical zu stellen. Bei \textsc{Reichenbach}'s Kreisen ist dazu eigentlich die
H"anglibelle bestimmt, die zwar grosse Genauigkeit geben kann, deren Gebrauch aber, weil das Fernrohr
abgenommen werden muss, etwas umst"andlich ist, eine sehr solide Aufstellung und sehr viel Behutsamkeit
erfordert. Ich bediene mich des Quecksilberhorizonts, w"ahle ein Object, so hoch, wie m"oglich (es braucht
"ubrigens nicht entfernt zu sein, da es bei dieser Operation nichts schadet, wenn man auch die Okularr"ohre,
um es deutlich zu sehen, weit herausziehen muss), bringe den nach dem Augenmaass vertical gestellten
Kreis in dessen Verticalebene, und pointire genau auf dasselbe. Hierauf wird das Fernrohr heruntergef"uhrt,
bis man das Bild des Objects im Quecksilberhorizont sieht. Kommt die Gesichtslinie genau auf dasselbe,
so steht die Kreisebene vertical; wo nicht, so corrigirt man die eine H"alfte entweder an einer Fussschraube,
oder an der Schraube ohne Ende, die andere an der Azimuthschraube. Hat man die Halbirung genau
getroffen, so wird die Gesichtslinie des wieder hinaufgef"uhrten Fernrohrs das Object selbst genau treffen,
sonst wird abermals die eine H"alfte der Abweichung an der Fussschraube oder der Schraube ohne Ende,
die andere an der Azimuthschraube corrigirt. Ist die Elevation des Objects nicht zu klein, und dieses
scharf genug zu pointiren, so l"asst sich auf diesem Wege auch eine sehr grosse Genauigkeit erreichen,
und eine viel gr"ossere, als n"othig ist, wenn man nicht sehr kleine Zenithdistanzen messen will. Auf gar
zu kleine Zenithdistanzen ist es aber bekanntlich nicht rathsam, den \textsc{Borda}ischen Kreis anzuwenden.

\newpage

\section*{INSTRUMENTE.}
\markboth{INSTRUMENTE.}{}
\addcontentsline{toc}{section}{Instrumente}

\begin{center}
\rule{0.6\textwidth}{0.4pt}\\[0.5em]
\textit{G"ottingische gelehrte Anzeigen. St"uck 127. Seite 1257\,..\,1267. 1818. August 8.}\\[0.3em]
\rule{0.6\textwidth}{0.4pt}
\end{center}

Von unserer \textit{neuen Sternwarte}, diesem grossen und sprechenden Denkmal der Liebe unserer Regierung
f"ur die Wissenschaft, ist bisher in unsern Bl"attern noch keine besondere Erw"ahnung geschehen, obgleich
das Geb"aude bereits seit anderthalb Jahren in dem Maasse vollendet ist, dass der Instrumentenvorrath der
alten Sternwarte in dasselbe aufgenommen, und unterbrochen in einer Abtheilung des Geb"audes beobachtet
werden konnte. Allein der Natur der Sache nach konnten diese Beobachtungen nur zu der Gattung derjenigen
geh"oren, dergleichen auch auf der alten Sternwarte schon sich anstellen liessen, und es schien uns nicht
passend, davon in diesen Bl"attern besondere Anzeige zu machen; verschiedene davon sind bereits in
astronomischen Zeitschriften bekannt gemacht. Die neue Sternwarte, bestimmt, keiner nachzustehen,
kann in den ihr geb"uhrenden Rang erst durch den Besitz der \textit{festen Meridian-Instrumente} treten, und
die Zeit ist jetzt nahe, wo sie vollst"andig ausger"ustet sein wird. --- Ueber das Aeussere des Geb"audes,
welches der W"urde seiner Bestimmung entspricht und der geschickten und geschmackvollen Ausf"uhrung
unsers Hrn.~Universit"ats-Baumeisters \textsc{M"ullee} zur Ehre gereicht, werden wir hier nichts sagen. Auch
eine vollst"andige Beschreibung der innern Einrichtung wird einem andern Orte vorbehalten bleiben.
Aber die neuen Hauptinstrumente, welche dem Geb"aude erst seinen wahren wissenschaftlichen Werth geben,
sollen, so wie sie nach und nach ankommen und aufgestellt werden, durch diese Bl"atter n"aher angezeigt,
und in die Bekanntschaft der Verehrer der Himmelskunde eingef"uhrt werden. --- Das erste der neuen
Meridianinstrumente, der \textit{Repsoldsche Meridiankreis}, kam im April d.~J. an und wurde von dem K"unstler
selbst aufgestellt. Dieses bereits vor l"angerer Zeit von Hrn.~\textsc{Repsold} in Hamburg urspr"unglich zu seinem
eignen Gebrauch verfertigte, und in dessen Privatsternwarte auf dem Hamburger Walle aufgestellt gewesene
Instrument, ist den Freunden der Astronomie nicht unbekannt, indem verschiedene Ausz"uge aus Hrn.~\textsc{Repsold}s
Tageb"uchern in der Monatlichen Correspondenz abgedruckt sind; auch in unsern Bl"attern (1811
S.~1290 [S.~327 d.~B.]) haben wir einige von Hrn.~Professor \textsc{Schumacher} an diesem Kreise gemachte
Beobachtungen angezeigt. Da Hrn.~\textsc{Repsold}s Sternwarte in der f"ur Hamburg so ungl"ucklichen Periode
demolirt war, und der Besitzer demnach das Instrument unbenutzt lassen musste, so genehmigte unsere
Regierung den Vorschlag des Hrn.~Hofrath \textsc{Gauss}, dasselbe f"ur unsere neue Sternwarte anzukaufen.
Der K"unstler "ubernahm dabei mehrere so wichtige und bedeutende Vervollkommnungen an dem Instrumente,
dass dieses gewissermassen ein ganz neues geworden ist. Die durch manche unvorhergesehene Umst"ande
oft unterbrochene Arbeit an diesen neuen Einrichtungen ist die Ursache, dass die Ablieferung des Instruments
sich so lange verz"ogert hat, zu dessen Aufstellung auf der Sternwarte bereits seit anderthalb Jahren
alles vorbereitet war. --- Die Einrichtung dieses Instruments verdient um so mehr eine umst"andlichere
Beschreibung, da sie, ganz aus den eignen Ideen des genialen K"unstlers hervorgegangen, bis jetzt einzig
ist. Es vereinigt in sich ein vollkommenes \textit{Mittagsfernrohr} mit einem Kreise, der dem neuen Greenwicher
Mauerkreise sehr "ahnlich ist.

Das Fernrohr von $7\frac{1}{2}$~Pariser Fuss L"ange, $6\frac{1}{2}$~Fuss Brennweite und 46~Linien Oeffnung, befindet
sich an einer Axe von 4~Fuss L"ange, die wie bei andern Mittagsfernr"ohren zwischen zwei steinernen
Pfeilern von $6\frac{1}{2}$~Fuss H"ohe und 22~Zoll im Quadrat, aufgeh"angt ist. Durch drei besondere Balancirungen
wird sowohl die Biegung des Fernrohrs, als die der Axe aufgehoben und das Gewicht des Ganzen so weit
getragen, dass die Pfannen nur einen Druck von einigen Lothen erleiden. Die Pfannen selbst sind von
Bergkrystall, die vollkommen cylindrischen und gleich dicken Zapfen von Glockenmetall; an letztem kann,
wenn sie nach langem Gebrauch durch die geringe Reibung doch etwas abgenutzt werden sollten, durch
eine besondere Einrichtung eine andere Stelle zum Aufliegen gebracht werden. Die Ocularr"ohre ist parallel
mit der Axe verschiebbar, um jeden Faden in die Mitte des Gesichtsfeldes bringen zu k"onnen; die Beleuchtung
der F"aden geht durch die Axe. Die st"arkste Vergr"osserung, welche gew"ohnlich gebraucht wird, ist eine
90malige. Bei g"unstiger Luft lassen sich Sterne bis zur dritten Gr"osse bei Tage ohne M"uhe beobachten.
Die H"anglibelle, wodurch die Axe horizontal gestellt wird, ist h"ochst sorgf"altig gearbeitet, und l"asst
Theile von Secunden mit Sicherheit erkennen.

Den beiden Pfeilern dient Eine grosse Platte zur gemeinschaftlichen Unterlage, welche selbst durch
grosse Quadern 12~Fuss tief im Boden fundirt ist. Die Solidit"at dieser Einrichtung hat sich seit der
Aufstellung auf das vollkommenste bew"ahrt; Azimuth, Horizontalit"at und Collimationslinie zeigen bis
jetzt wenigstens eine fast absolute Unwandelbarkeit. --- Auf der Axe, nahe dem einen Ende, sitzt der Kreis
fest, der bei der Theilung $3\frac{1}{2}$~Fuss im Durchmesser h"alt. Die Theilung ist mit einer bewundernsw"urdigen
Genauigkeit durch sehr saubere Striche von 5 zu 5 Minuten gemacht, w"ahrend der Kreis auf der Axe
selbst sass, wodurch alle Excentricit"at vermieden ist. Zur Versendung musste er freilich wieder abgenommen
werden; durch die an Ort und Stelle geschehene Wiederaufsetzung ist nur eine fast unmerkliche
Excentricit"at entstanden, die nach der sorgf"altigen Pr"ufung des Hrn.~Hofr.~\textsc{Gauss} $0''8$ betr"agt.
Die Ablesung geschieht durch drei vortreffliche mikrometrische Mikroskope, deren Tr"ager an den schiebbaren
Theilen der Lager so befestigt sind, dass ihre Entfernung von der Kreisfl"ache, auch wenn diese von der
Ebene etwas abweichen sollte, vollkommen best"andig bleibt; das eine Mikroskop sitzt unten, die beiden
andern seitw"arts $90^{\circ}$ von jenem entfernt. Ein viertes festes Mikroskop anzubringen verstattete der Bau
des Instruments nicht (fr"uher war nur ein einziges da). Da die Wirkung einer Excentricit"at durch die
beiden Seitenmikroskope ganz aufgehoben wird, so bleibt bei der Ablesung aller, nur der dritte Theil
davon "ubrig, der nur auf $0''3$ sich belaufen, und wo es n"othig scheint mit in Rechnung gebracht werden
kann. Ein viertes bewegliches Mikroskop zur Pr"ufung der Theilung selbst, wird der K"unstler auf den
Wunsch des Hrn.~Hofr.~\textsc{Gauss} noch nachliefern. Die L"aufer in den Mikroskopen, welche durch "ausserst
gleichf"ormige von allem todten Gange durch die bekannte \textsc{Ramsden}sche Erfindung frei gemachte Schrauben
bewegt werden, f"uhren ein kreisrundes Loch, in welchem die Bissection der Theilstriche mit einer solchen
Sch"arfe erkannt werden kann, dass man auf Theile von Secunden sicher ist. Der Kreis hat neun Speichen,
zwischen einem Paar derselben befindet sich ein kupferner Cylinder, welcher durch eine ebenfalls h"ochst
vortreffliche Libelle auf eine halbe Secunde genau horizontal gestellt werden kann. Hierdurch werden nicht
allein die etwanigen Ver"anderungen in der Stellung der Mikroskope bemerkbar, sondern, da die Neigung
der Gesichtslinie gegen diesen Cylinder durch Beobachtungen in gewechselter Lage des Instruments bekannt
wird, so erh"alt man auch absolute Zenithdistanzen. --- Bei der Ankunft des Instruments war das Fadennetz
im Fernrohr noch dasselbe, welches bei der ersten Verfertigung eingezogen war, und aus drei Verticalet,
und zwei horizontalen $1''$ von einander abstehenden Spinnenf"aden bestand. Da Hr.~\textsc{Repsold} Ende
Mai wiederum nach G"ottingen zur"uckkam, so vertauschte er dieses Fadennetz mit einem neuen, wodurch
das Pointiren an Genauigkeit betr"achtlich gewonnen hat. Das neue Fadennetz besteht aus f"unf verticalen
und zwei horizontalen ausgesuchten Spinnenf"aden; letztere sind vollkommen parallel, und stehen $1''5$
von einander ab. Beim Beobachten wird in der Regel der Stern in die Mitte zwischen diese beiden
F"aden gebracht, welches mit sehr grosser Sch"arfe geschehen, und wobei man die Antritte an die verticalen
F"aden immer leicht und scharf beobachten kann. Die Zwischenzeiten zwischen den einzelnen Verticalf"aden
betragen f"ur Sterne im Aequator etwas "uber 16~Secunden. --- Die Beobachtungen des Polarsterns
geh"oren auf jeder wohlbestellten Sternwarte aus bekannten Gr"unden zu den t"aglichen; bei einem neuen
Mittagsfernrohr sind sie aber doppelt wichtig, da auf dieselben die Hauptberichtigungen gegr"undet werden
m"ussen. Hr.~Hofr.~\textsc{Gauss} hat daher, so oft es das im Juni sehr g"unstige Wetter erlaubte, alle obern
und untern Culminationen dieses Sterns sorgf"altig beobachtet. Wir geben hier als eine Probe von dem,
was das Instrument leistet, die bisher beobachteten geraden Aufsteigungen, nebst der Vergleichung mit der
von Hrn.~\textsc{Bessel} im astronomischen Jahrbuche 1817 gegebenen und im folgenden Jahrgange verbesserten
H"ulfstafel.

\newpage

Der Bruch des Tages ist von der untern Culmination, welche zu dem angesetzten Datum geh"ort,
gez"ahlt. Die letzte Columne dr"uckt das Quadrat der Genauigkeit jedes einzelnen Resultats aus, nach
Massgabe der Anzahl der concurrirenden Fadenantritte, welches Quadrat von einigen Astronomen ganz
schicklich das \textit{Gewicht} der Beobachtung genannt ist; als Einheit liegt dabei diejenige Genauigkeit zum
Grunde, die eine aus zwei einander folgenden Culminationen, die jede nur an Einem Faden beobachtet
sind, geschlossene Bestimmung hat. Die folgenden 19 Bestimmungen gr"undeten sich auf 41 beobachtete
Culminationen, von denen immer zwei oder drei unmittelbar auf einander folgende zu Einem Resultate
verbunden wurden.

\begin{center}
\small
\begin{tabular}{c|c|c|c|c}
\toprule
1818 & Beob. & gerade Aufsteigung & Correction & Gewicht \\
& & des Nordsterns & der Tafel & \\
\midrule
Juni \phantom{0}3.25 & & $0^{\mathrm{h}}\;55^{\mathrm{m}}\;58^{\mathrm{s}}65$ & & $1.67$ \\
\phantom{0}\phantom{0}\phantom{0} 6.25 & & & & \\
\phantom{0}\phantom{0}\phantom{0} 4.25 & & $\phantom{0}\phantom{0}\phantom{0}\phantom{0}\phantom{0}\phantom{0}\phantom{0}58.19$ & $+\,3''93$ & $2.45$ \\
\phantom{0}\phantom{0}\phantom{0} 5.25 & & & $+\,2.80$ & \\
\phantom{0}\phantom{0}\phantom{0} 8.50 & & & $+\,2.97$ & $6.40$ \\
\phantom{0}\phantom{0} 10.25 & & & & $3.75$ \\
\phantom{0}\phantom{0} 11.25 & & $0\;56\;58.56$ & & \\
& & $\phantom{0}\phantom{0}\phantom{0}\phantom{0}\phantom{0}\phantom{0}\phantom{0}\phantom{0}1.14$ & & \\
& & $\phantom{0}\phantom{0}\phantom{0}\phantom{0}\phantom{0}\phantom{0}\phantom{0}58.91$ & $+\,2.22$ & \\
\phantom{0}\phantom{0} 16.75 & & & $+\,3.20$ & $4.29$ \\
\phantom{0}\phantom{0} 12.50 & & $\phantom{0}\phantom{0}\phantom{0}\phantom{0}\phantom{0}\phantom{0}\phantom{0}\phantom{0}2.58$ & & $1.33$ \\
& & & $+\,2.59$ & \\
& & $\phantom{0}\phantom{0}\phantom{0}\phantom{0}\phantom{0}\phantom{0}\phantom{0}\phantom{0}2.68$ & & \\
& & $\phantom{0}\phantom{0}\phantom{0}\phantom{0}\phantom{0}\phantom{0}\phantom{0}\phantom{0}8.29$ & & \\
\phantom{0}\phantom{0} 18.25 & & & & $4.44$ \\
& & $\phantom{0}\phantom{0}\phantom{0}\phantom{0}\phantom{0}\phantom{0}\phantom{0}10.49$ & $+\,2.36$ & $5.00$ \\
\phantom{0}\phantom{0} 19.25 & & $\phantom{0}\phantom{0}\phantom{0}\phantom{0}\phantom{0}\phantom{0}\phantom{0}\phantom{0}6.11$ & $+\,2.13$ & \\
\phantom{0}\phantom{0} 21.25 & & & & \\
\phantom{0}\phantom{0} 22.75 & & $\phantom{0}\phantom{0}\phantom{0}\phantom{0}\phantom{0}\phantom{0}\phantom{0}\phantom{0}3.31$ & & \\
\phantom{0}\phantom{0} 26.25 & & $\phantom{0}\phantom{0}\phantom{0}\phantom{0}\phantom{0}\phantom{0}\phantom{0}13.49$ & $+\,2.51$ & $1.67$ \\
& & $\phantom{0}\phantom{0}\phantom{0}\phantom{0}\phantom{0}\phantom{0}\phantom{0}14.77$ & $+\,2.26$ & $2.86$ \\
\phantom{0}\phantom{0} 27.25 & & & & \\
\phantom{0}\phantom{0}\phantom{0} 8.25 & & $\phantom{0}\phantom{0}\phantom{0}\phantom{0}\phantom{0}\phantom{0}\phantom{0}\phantom{0}7.31$ & & $4.44$ \\
& & $\phantom{0}\phantom{0}\phantom{0}\phantom{0}\phantom{0}\phantom{0}\phantom{0}12.02$ & $+\,3.26$ & $3.43$ \\
\phantom{0}\phantom{0} 29.50 & & & $+\,2.60$ & \\
& & & $+\,3.60$ & \\
Juli \phantom{0}\phantom{0} 7.50 & & $\phantom{0}\phantom{0}\phantom{0}\phantom{0}\phantom{0}\phantom{0}\phantom{0}22.51$ & $+\,3.14$ & $5.00$ \\
\phantom{0}\phantom{0}\phantom{0} 9.25 & & $\phantom{0}\phantom{0}\phantom{0}\phantom{0}\phantom{0}\phantom{0}\phantom{0}22.97$ & & $3.75$ \\
Sept. \phantom{0} 3.75 & & $\phantom{0}\phantom{0}\phantom{0}\phantom{0}\phantom{0}\phantom{0}\phantom{0}23.89$ & $+\,3.42$ & $6.00$ \\
& & $\phantom{0}\phantom{0}\phantom{0}\phantom{0}\phantom{0}\phantom{0}\phantom{0}57.44$ & $+\,2.35$ & $5.00$ \\
& & & $+\,3.31$ & \\
& & & $+\,3.22$ & $1.60$ \\
& & & $+\,3.40$ & \\
\bottomrule
\end{tabular}
\end{center}

Das Mittel der durch diese 19 Beobachtungen gefundenen Correctionen der Tafel ist $+\,2''88$. Insofern
man sich erlaubt, den Tafelfehler w"ahrend dieser Zeit als best"andig zu betrachten, gibt die Anwendung
der Wahrscheinlichkeitstheorie des Hrn.~Hofrath \textsc{Gauss} den wahrscheinlichen Fehler einer Beobachtung,
deren Genauigkeit 1 ist, $= 0''64$, und den wahrscheinlichen Fehler des Endresultats $= 0''076$. Wir f"ugen
ferner, als eine zweite Probe, die Beobachtungen des \textsc{Uranus} bei, welche Hr.~Hofr.~\textsc{Gauss} zur Zeit der
Opposition dieses Planeten angestellt hat. Die Declinationen gr"unden sich auf die Vergleichung mit der
unmittelbar vorausgehenden untern und der unmittelbar folgenden obern Culmination des Nordsterns, blos
die vom 7ten und 13ten Juni ausgenommen, wo blos die untere Culmination beobachtet war, und wo die
Declination des Nordsterns aus \textsc{Bessel}s Tafel zum Grunde gelegt ist. Die Refractionen sind aus
\textsc{Bessel}s Tafel, die sich auf die \textsc{Bradley}schen Beobachtungen gr"undet, genommen; w"urden dieselben
aus \textsc{Carlini}s Tafel entlehnt, so w"urden die Declinationen des \textsc{Uranus} $3''$ s"udlicher ausfallen.
Wenn in Zukunft erst eine hinl"angliche Anzahl schicklicher Beobachtungen beisammen sein wird, wird sich
zeigen, welche von beiden Tafeln dem hiesigen Clima am angemessensten ist.

\newpage

\begin{center}
\textbf{Beobachtungen des Uranus.}
\end{center}

\begin{center}
\small
\begin{tabular}{c|c|c|c|c}
\toprule
1818 & Mittlere Zeit & Gerade Aufsteig. & \multicolumn{2}{c}{S"udl.~Abweich.} \\
\midrule
Juni \phantom{0}3 & $12^{\mathrm{h}}\;1^{\mathrm{m}}\;22^{\mathrm{s}}17^{\mathrm{s}}5$ & & & $1^{\circ}\;1'\;41''5$ \\
& $12\;\phantom{0}3\;18\;10.9$ & $257^{\circ}\;20'\;\phantom{0}\phantom{0}9.98$ & & \phantom{0}\phantom{0}\phantom{0}$\phantom{0}29.0$ \\
\phantom{0}\phantom{0} 4 & $12\;14\;\phantom{0}4.6$ & \phantom{0}\phantom{0}\phantom{0}$23'\;28''3$ & & \phantom{0}\phantom{0}\phantom{0}$\phantom{0}50.5$ \\
\phantom{0}\phantom{0} 5 & $12\;12\;22$ & \phantom{0}\phantom{0}\phantom{0}\phantom{0}$\phantom{0}\phantom{0}\phantom{0}54.6$ & & \phantom{0}\phantom{0}\phantom{0}$\phantom{0}38.4$ \\
\phantom{0}\phantom{0} 7 & $12\;\phantom{0}5\;51.6$ & & & \phantom{0}\phantom{0}\phantom{0}$\phantom{0}27.4$ \\
\phantom{0}\phantom{0} 8 & $12\;\phantom{0}1\;45.1$ & \phantom{0}\phantom{0}$\phantom{0}\phantom{0}10\;15.0$ & & \phantom{0}\phantom{0}\phantom{0}$\phantom{0}\phantom{0}2.9$ \\
\phantom{0}\phantom{0} 9 & $11\;57\;38.6$ & & & \phantom{0}\phantom{0}\phantom{0}$\phantom{0}12.6$ \\
\phantom{0}\phantom{0} 10 & $11\;53\;32.1$ & $27\;36.4$ & & \\
\phantom{0}\phantom{0} 11 & $11\;49\;25.7$ & \phantom{0}\phantom{0}$\phantom{0}\phantom{0}\phantom{0}18.0$ & $22^{\circ}\;59'\;50.9$ & \\
\phantom{0}\phantom{0} 12 & $11\;45\;19.2$ & $256^{\circ}\;59'\;56.5$ & & \phantom{0}\phantom{0}$\phantom{0}\phantom{0}6.2$~S"udl. \\
& $11\;41\;12.7$ & \phantom{0}\phantom{0}\phantom{0}$\phantom{0}\phantom{0}33.9$ & \phantom{0}\phantom{0}$\phantom{0}59\;40.5$ & \\
\bottomrule
\end{tabular}
\end{center}

Hr.~Hofr.~\textsc{Gauss} "ubertrug die Vergleichung dieser Beobachtungen mit den \textsc{Delambre}schen
\textsc{Uranus}-Tafeln Hrn.~\textsc{Dirksen}, welcher sich bei uns dem Studium der Astronomie mit ausgezeichnetem
Eifer widmet. Die Resultate dieser Vergleichung sind folgende:

\begin{center}
\small
\begin{tabular}{c|c|c|c}
\toprule
& \multicolumn{2}{c|}{Unterschiede} & \\
& Ger.~Aufst. & Abweichung & \\
\midrule
Juni \phantom{0}3 & $-\,54.1$ & & \\
\phantom{0}\phantom{0} 4 & $-\,52.9$ & $+\,10.3$ & \\
\phantom{0}\phantom{0} 5 & $-\,52.1$ & & \\
\phantom{0}\phantom{0} 7 & $-\,55.8$ & $+\,12.9$ & \\
\phantom{0}\phantom{0} 8 & $-\,55.9$ & $+\,10.6$ & \\
\phantom{0}\phantom{0} 9 & $-\,56.4$ & & \\
\phantom{0}\phantom{0} 10 & $-\,56.8$ & $+\,12.3$ & \\
\phantom{0}\phantom{0} 11 & $-\,57.2$ & $+\,10.5$ & \\
\phantom{0}\phantom{0} 12 & $-\,52.4$ & $+\,\phantom{0}9.4$ & \\
\midrule
& & $+\,12.9$ & \\
& & $+\,12.0$ & \\
& $-\,58.1$ & & \\
& $-\,55.2$ & $+\,10.1$ & \\
\midrule
Mittel & & $+\,10.8$ & \\
& & $+\,\phantom{0}7.4$ & \\
\bottomrule
\end{tabular}
\end{center}

Hieraus ergab sich, nach Hrn.~\textsc{Dirksen}s Rechnung:

\begin{center}
\begin{tabular}{l@{\quad}l}
Zeit der Opposition 1818 Juni 9. & $5^{\mathrm{h}}\;30^{\mathrm{m}}\;43^{\mathrm{s}}$ Mittlere Zeit in G"ottingen \\
Wahre L"ange \dotfill & $258^{\circ}\;10'\;34''5$ \\
Geocentrische Breite \dotfill & $\phantom{0}0^{\circ}\;\phantom{0}4'\;12.1$ \\
Heliocentrische \dotfill & $\phantom{0}0\;\phantom{0}3\;58.7$~S"udl.
\end{tabular}
\end{center}

und der Fehler der \textsc{Delambre}schen Tafeln
\begin{center}
\begin{tabular}{l@{\dotfill}l}
in der L"ange & $-\,45''3$ \\
in Heliocentrischer Breite & $+\,14.6$ \\
\end{tabular}
\end{center}

Wir k"onnen diese Anzeige nicht schliessen, ohne zugleich der \textsc{Shelton}schen Pendeluhr zu erw"ahnen,
die zwar schon seit 47 Jahren auf der alten Sternwarte im Gebrauch gewesen war, deren Vortrefflichkeit
aber erst jetzt ganz gew"urdigt werden kann, da fr"uher die Sternwarte kein Mittel zur allersch"arfsten
Zeitbestimmung besass. Wir f"ugen zum Beweise nur das Register ihres Ganges von einem Monat bei,
und bemerken, dass dieselbe Gleichf"ormigkeit w"ahrend der ganzen Zeit statt gefunden hat, wo das
\textsc{Repsold}sche Instrument zur Bestimmung des Ganges gedient hat. Das Datum bezieht sich auf Sonnentage,
und die Stunden der Uhrzeit sind, um die Uebersicht zu erleichtern jedesmal "uber 24 hinaus bis zum
n"achsten Mittage fortgez"ahlt; die letzte Columne enth"alt die Anzahl der Sterne des \textsc{Maskelyne}schen
Fundamentalcatalogs, auf welche die Zeitbestimmung gegr"undet ist.

\newpage

\begin{center}
\small
\begin{tabular}{c|c|c|c|c|c|c}
\toprule
& Uhrzeit & Stand gegen & Verglichene & & Uhrzeit & Stand gegen \\
& & Sternzeit & Sterne & & & Sternzeit \\
\midrule
Juni 15 & $20^{\mathrm{h}}\;50$ & $+0^{\mathrm{s}}41$ & 4 & Juli \phantom{0}1 & $14^{\mathrm{h}}\;51$ & $-\,3^{\mathrm{s}}11$ \\
\phantom{15} 23 & $11\;37$ & & 2 & \phantom{0}\phantom{0}\phantom{0} 2 & $25\;17$ & \\
\phantom{15} 18 & $15^{\mathrm{h}}\;22$ & $+\,0.29$ & & \phantom{0}\phantom{0}\phantom{0} 6 & $22\;31$ & $-\,4.53$ \\
& & $+\,0.28$ & 2 & \phantom{0}\phantom{0}\phantom{0} 7 & $28^{\mathrm{h}}\;14$ & \\
\phantom{15} 27 & $12\;28$ & $-\,0.06$ & 2 & \phantom{0}\phantom{0}\phantom{0} 8 & $15\;\phantom{0}3$ & $-\,4.28$ \\
\phantom{15} 20 & & & 2 & \phantom{0}\phantom{0} 10 & & $-\,4.93$ \\
\phantom{15} 21 & $29$ & $+\,0.07$ & & \phantom{0}\phantom{0}\phantom{0} 9 & & $-\,4.62$ \\
\phantom{15} 22 & & $-\,1.09$ & 1 & & & $-\,4.73$ \\
\phantom{15} & $11\;40$ & & 1 & \phantom{0}\phantom{0} 16 & $15\;56$ & \\
\phantom{15} 26 & $16\;15$ & $-\,1.86$ & 5 & \phantom{0}\phantom{0} 29 & $14\;41$ & \\
\phantom{15} 28 & $14\;41$ & $-\,2.43$ & 4 & \phantom{0}\phantom{0} 15 & & \\
& $18$ & $-\,2.32$ & 1 & & & \\
\bottomrule
\end{tabular}
\end{center}

Hiebei darf nicht unerw"ahnt bleiben, dass die Aufstellung der Uhr erst noch provisorisch ist.
Sie ruht zwar auf einem besondern steinernen Fundament, aber nur vermittelst ihres Geh"auses, ohne
unmittelbar an den steinernen Pfeiler, neben welchen dieses gestellt ist, befestigt zu sein. Sie geht einen
Monat in Einem Aufzuge; um die Mitte dieser Zeit kommt das Gewicht der Linse gegen"uber zu stehen,
welches man sonst wohl f"ur nachtheilig gehalten und durch mancherlei Einrichtungen zu vermeiden
gesucht hat. Allein das Register der \textsc{Shelton}schen Uhr zeigt um diese Zeit gar keine sp"urbare Ver"anderung
im Gange, der "uberhaupt viel regelm"assiger ist, als der mancher andern Uhren auf den ersten Sternwarten,
auch solcher, die man unmittelbar an einen steinernen Pfeiler befestigt hat. Sollte man nicht hieraus
schliessen, dass beide Umst"ande von viel geringerer Wichtigkeit sind, als man bisher gew"ohnlich geglaubt hat?

\newpage

\subsection*{GAUSS AN BODE.}
\markboth{GAUSS AN BODE.}{}

\begin{center}
\rule{0.6\textwidth}{0.4pt}\\[0.5em]
\textit{Astronomisches Jahrbuch f"ur 1821. Seite 212\,..\,216. Berlin 1818.}\\[0.3em]
\rule{0.6\textwidth}{0.4pt}
\end{center}

\begin{flushright}
G"ottingen, den 7.~Sept.~1818.
\end{flushright}

Ihrem Wunsche gem"ass "ubersende ich Ihnen, hochverehrter Freund, beigehend abermals eine
Ephemeride f"ur die Pallas vom August 1819 bis Juni 1820. Sie ist vom Hrn.~\textsc{Dirksen} berechnet, welcher
sich hier mit Eifer den mathematischen Wissenschaften widmet. Eben derselbe hat auch die Berechnung
der Opposition des laufenden Jahres noch einmal wiederholt nach denselben Elementen, nach welchen
Hr.~Dr.~\textsc{Tittel} gerechnet hatte, und ein etwas abweichendes Resultat gefunden, n"amlich:

\begin{center}
\begin{tabular}{l|c|c}
& Hr.~\textsc{Dirksen} & Hr.~Dr.~\textsc{Tittel} \\
\midrule
Zeit in G"ottingen & Sept.~6\ $8^{\mathrm{h}}\;47^{\mathrm{m}}\;15^{\mathrm{s}}$ & Sept.~8\ $18^{\mathrm{h}}\;44^{\mathrm{m}}\;12^{\mathrm{s}}$ \\
Wahre L"ange & $345^{\circ}\;47'\;17''$ & $345^{\circ}\;55'\;34''$ \\
Geocentrische Breite & $\phantom{0}6\;47\;\phantom{0}5$ & $\phantom{0}6\;47\;11$ \\
\end{tabular}
\end{center}

Welches die Entstehung dieses Rechnungsfehlers sei, kann ich nicht sagen; er ist aber auf die
Berechnung der Opposition eingeschr"ankt, da mehrere von Hrn.~\textsc{Dirksen} nachgerechnete Oerter der
Ephemeride vollkommen mit Hrn.~Dr.~\textsc{Tittel}s Angaben "ubereinstimmen (deren Richtigkeit auch durch
meine nachher anzuf"uhrenden Beobachtungen bereits best"atigt wird).

Seit Ende Aprils d.~J. beobachte ich nunmehro an dem trefflichen Meridiankreise von \textsc{Repsold},
welches Instrument Sie ohne Zweifel fr"uher selbst gesehen haben, woran aber jetzt so viele wichtige
Verbesserungen angebracht sind, dass es f"ur ein ganz neues gelten kann. Es hat eine ganz neue Theilung
erhalten, drei Ablesungen statt der fr"uhern einzigen, neue Zapfenlager von Bergkrystall, die Zapfen selbst
sind neu abgedreht und centrirt u.s.w. Verschiedene andere neue Vorrichtungen dazu erwarte ich noch
von Hrn.~\textsc{Repsold}. Als Probe von der Vortrefflichkeit des Instruments, insofern es Mittagsfernrohr
ist, m"ogen folgende Bestimmungen der Ger.~Aufsteig.~des Nordsterns dienen; die letzte Columne unter
der Ueberschrift \textit{Gewicht} dr"ucken das Quadrat der Genauigkeit nach Massgabe der Zahl der concurrirenden
F"aden aus, wenn die Genauigkeit einer Bestimmung aus einem Fadenantritt in the obern und einem
Fadenantritt in der untern Culmination zur Einheit genommen wird. Die Formel daf"ur ist
\[\frac{8abc}{ab + 4ac + bc}\]
wenn, bei einer auf zwei Culminationen gegr"undeten Bestimmung $c$ und $b$ die Anzahl der beobachteten
Antritte in jeder ausdr"ucken, und
\[\frac{8abc}{ab + 4ac + bc}\]
bei einer Bestimmung aus drei auf einander folgenden Culminationen, wo $c$ F"adenantritte in der dritten
beobachtet sind. Im ersten Falle setzt man voraus, dass die Pfeiler 12 Stunden unverr"uckt gestanden;
im letztern nur, dass ihre etwanigen Ver"anderungen w"ahrend 24 Stunden gleichf"ormig gewesen sind.

\newpage

\begin{center}
\small
\begin{tabular}{c|c|c|c|c}
\toprule
1818 & Beob. & Ger.~Aufst. & Correction von & Gewicht \\
& & des Nordsterns & \textsc{Bessel}s Tafel & \\
\midrule
Juni \phantom{0}3.25 & & $0^{\mathrm{h}}\;55^{\mathrm{m}}\;58^{\mathrm{s}}65$ & & $1.67$ \\
\phantom{0}\phantom{0}\phantom{0} 6.25 & & & & \\
\phantom{0}\phantom{0}\phantom{0} 4.25 & & $\phantom{0}\phantom{0}\phantom{0}\phantom{0}\phantom{0}\phantom{0}\phantom{0}58.19$ & $++\,3^{\mathrm{s}}93$ & $2.45$ \\
\phantom{0}\phantom{0}\phantom{0} 5.25 & & & $+\,2.80$ & \\
\phantom{0}\phantom{0}\phantom{0} 8.50 & & & $+\,2.97$ & $6.40$ \\
\phantom{0}\phantom{0} 10.25 & & & & $3.75$ \\
\phantom{0}\phantom{0} 11.25 & & $0\;56\;58.56$ & & \\
& & $\phantom{0}\phantom{0}\phantom{0}\phantom{0}\phantom{0}\phantom{0}\phantom{0}\phantom{0}1.14$ & & \\
& & $\phantom{0}\phantom{0}\phantom{0}\phantom{0}\phantom{0}\phantom{0}\phantom{0}58.91$ & $+\,2.22$ & \\
\phantom{0}\phantom{0} 16.75 & & & $+\,3.20$ & $4.29$ \\
\phantom{0}\phantom{0} 12.50 & & $\phantom{0}\phantom{0}\phantom{0}\phantom{0}\phantom{0}\phantom{0}\phantom{0}\phantom{0}2.58$ & & $1.33$ \\
& & & $+\,2.59$ & \\
& & $\phantom{0}\phantom{0}\phantom{0}\phantom{0}\phantom{0}\phantom{0}\phantom{0}\phantom{0}2.68$ & & \\
& & $\phantom{0}\phantom{0}\phantom{0}\phantom{0}\phantom{0}\phantom{0}\phantom{0}\phantom{0}8.29$ & & \\
\phantom{0}\phantom{0} 18.25 & & & & $4.44$ \\
& & $\phantom{0}\phantom{0}\phantom{0}\phantom{0}\phantom{0}\phantom{0}\phantom{0}10.49$ & $+\,2.36$ & $5.00$ \\
\phantom{0}\phantom{0} 19.25 & & $\phantom{0}\phantom{0}\phantom{0}\phantom{0}\phantom{0}\phantom{0}\phantom{0}\phantom{0}6.11$ & $+\,2.13$ & \\
\phantom{0}\phantom{0} 21.25 & & & & \\
\phantom{0}\phantom{0} 22.75 & & $\phantom{0}\phantom{0}\phantom{0}\phantom{0}\phantom{0}\phantom{0}\phantom{0}\phantom{0}3.31$ & $+\,2.51$ & \\
\phantom{0}\phantom{0} 26.25 & & $\phantom{0}\phantom{0}\phantom{0}\phantom{0}\phantom{0}\phantom{0}\phantom{0}13.49$ & $+\,2.26$ & $1.67$ \\
& & $\phantom{0}\phantom{0}\phantom{0}\phantom{0}\phantom{0}\phantom{0}\phantom{0}14.77$ & & $2.86$ \\
\phantom{0}\phantom{0} 27.25 & & & & \\
\phantom{0}\phantom{0}\phantom{0} 8.25 & & $\phantom{0}\phantom{0}\phantom{0}\phantom{0}\phantom{0}\phantom{0}\phantom{0}\phantom{0}7.31$ & $+\,3.26$ & $4.44$ \\
& & $\phantom{0}\phantom{0}\phantom{0}\phantom{0}\phantom{0}\phantom{0}\phantom{0}12.02$ & $+\,2.60$ & $3.43$ \\
\phantom{0}\phantom{0} 29.50 & & & $+\,3.60$ & \\
& & & $+\,3.14$ & \\
Juli \phantom{0}\phantom{0} 7.50 & & $\phantom{0}\phantom{0}\phantom{0}\phantom{0}\phantom{0}\phantom{0}\phantom{0}22.51$ & & $5.00$ \\
\phantom{0}\phantom{0}\phantom{0} 9.25 & & $\phantom{0}\phantom{0}\phantom{0}\phantom{0}\phantom{0}\phantom{0}\phantom{0}22.97$ & $+\,3.42$ & $3.75$ \\
Sept. \phantom{0} 3.75 & & $\phantom{0}\phantom{0}\phantom{0}\phantom{0}\phantom{0}\phantom{0}\phantom{0}23.89$ & $+\,2.35$ & $6.00$ \\
& & $\phantom{0}\phantom{0}\phantom{0}\phantom{0}\phantom{0}\phantom{0}\phantom{0}57.44$ & $+\,3.31$ & $5.00$ \\
& & & $+\,3.22$ & \\
& & & $+\,3.40$ & $1.60$ \\
\bottomrule
\end{tabular}
\end{center}

In sofern man die Correction von \textsc{Bessel}s Tafel w"ahrend dieser Zeit als constant annimmt, w"are
dieselbe im Mittel
\[= +2^{\mathrm{s}}86.\]

\newpage

\section*{COMET.}
\markboth{COMET.}{}

\begin{center}
\rule{0.6\textwidth}{0.4pt}\\[0.5em]
\textit{G"ottingische gelehrte Anzeigen. St"uck 28. Seite 273\,..\,278. 1819 Febr.~18.}\\[0.3em]
\rule{0.6\textwidth}{0.4pt}
\end{center}

Ueber die beiden von \textsc{Pons} in Marseille im November des vorigen Jahrs entdeckten Cometen
theilen wir hier einige Nachrichten um so lieber mit, da diese Cometen in mehreren Beziehungen zu den
vorz"uglich merkw"urdigen geh"oren. Der eine scheint in Deutschland zuerst auf der Mannheimer Sternwarte
beobachtet zu sein. Die Beobachtungen, welche daselbst von Hrn.~Prof.~\textsc{Nicolai} angestellt und in einem
Schreiben an den Hrn.~H"ofrath \textsc{Gauss} vom 8.~Januar mitgetheilt sind, waren folgende:

\begin{center}
\small
\begin{tabular}{c|c|c|c|c}
\toprule
1818 & Mittlere Zeit & & Gerade & Nordliche \\
& in Mannheim & & Aufsteigung & Abweichung \\
\midrule
Dec. 22 & $5^{\mathrm{h}}\;45^{\mathrm{m}}\;37^{\mathrm{s}}$ & & $326^{\circ}\;16'\;\phantom{0}\phantom{0}\phantom{0}$ & $2^{\circ}\;54'\;\phantom{0}\phantom{0}5''$ \\
\phantom{22} 25 & $6\;51\;48$ & & $326\;\phantom{0}3\;59$ & $2\;52\;24$ \\
& & & & \phantom{0}\phantom{0}\phantom{0}$\phantom{0}22\;26$ \\
\phantom{22} 27 & $7\;14\;48$ & & $325\;48\;33$ & $2\;13\;34$ \\
\phantom{22} 29 & $7\;\phantom{0}0\;31$ & & $325\;\phantom{0}3\;49$ & $1\;10\;57$ \\
\phantom{22} & & & $324\;19\;58$ & \\
\bottomrule
\end{tabular}
\end{center}

Zugleich waren folgende parabolische Elemente beigef"ugt, welche Hr.~Prof.~\textsc{Nicolai} vorl"aufig aus
den Beobachtungen vom 22., 25.~und 29.~December berechnet hatte.

\begin{center}
\begin{tabular}{l@{\dotfill}l}
Zeit der Sonnenn"ahe 1819 Jan. & $25.0000$ in M. \\
Logarithm des kleinsten Abstandes & $0.52933$ \\
L"ange der Sonnenn"ahe & $145^{\circ}\;49'\;10''$ \\
L"ange des aufsteigenden Knoten & $330\;14\;17$ \\
Neigung der Bahn & $\phantom{0}14\;59\;\phantom{0}6$ \\
Richtung: rechtl"aufig. & \\
\end{tabular}
\end{center}

Auf der hiesigen Sternwarte wurden durch Hrn.~Prof.~\textsc{Harding} folgende Beobachtungen gemacht:

\begin{center}
\small
\begin{tabular}{c|c|c|c}
\toprule
& Mittlere Zeit & Gerade & Abweichung \\
& in G"ottingen & Aufsteigung & \\
\midrule
1818 Dec. 25 & $6^{\mathrm{h}}\;39^{\mathrm{m}}\;52^{\mathrm{s}}$ & $325^{\circ}\;33'\;\phantom{0}\phantom{0}$ & $2^{\circ}\;14'\;49''$~N"ordl. \\
\phantom{25} 26 & $5\;56\;58$ & $325\;17\;30$ & $2\;\phantom{0}0\;31$ \\
\phantom{25} 28 & $6\;24\;\phantom{0}8$ & $324\;59\;15$ & $1\;27\;31$ \\
1819 Jan. \phantom{0}1 & $6\;36\;16$ & $324\;40\;17$ & $1\;45\;\phantom{0}\phantom{0}$ \\
\phantom{0}\phantom{0} \phantom{0}7 & $6\;39\;41$ & $323\;11\;48$ & $0\;14\;58$ \\
\phantom{0}\phantom{0} 12 & $7\;\phantom{0}6\;\phantom{0}7$ & $319\;53\;59$ & \phantom{0}\phantom{0}$22\;51$ \\
\phantom{0}\phantom{0} & $7\;37\;34$ & $319\;15\;28$ & \phantom{0}$\phantom{0}19\;49$~S"udl. \\
\phantom{0}\phantom{0} & & $315\;36\;53$ & \phantom{0}$\phantom{0}\phantom{0}5\;35\;24$ \\
\bottomrule
\end{tabular}
\end{center}

Auf der Seeberger Sternwarte beobachtete Hr.~Prof.~\textsc{Encke} diesen Cometen f"unfmal;

\begin{center}
\small
\begin{tabular}{c|c|c|c}
\toprule
1819 & Mittlere Zeit & Gerade & \\
& in Seeberg & Aufsteigung & Abweichung \\
\midrule
Jan. \phantom{0}1 & $6^{\mathrm{h}}\;38^{\mathrm{m}}\;46^{\mathrm{s}}$ & $323^{\circ}\;11'\;36''$ & $1^{\circ}\;0'\;38''$~N"ordl. \\
\phantom{0}\phantom{0} \phantom{0}4 & $7\;26\;56$ & $321\;43\;43$ & $1\;0\;54\;59$~S"udl. \\
\phantom{0}\phantom{0} \phantom{0}5 & $7\;12\;12$ & $321\;10\;\phantom{0}3$ & $1\;20\;57$ \\
\phantom{0}\phantom{0} \phantom{0}6 & $6\;16\;32$ & $320\;35\;14$ & $1\;47\;57$ \\
\phantom{0}\phantom{0} & $5\;49\;48$ & $315\;35\;20$ & $5\;35\;26$ \\
\bottomrule
\end{tabular}
\end{center}

\newpage

Die Resultate welche Hr.~Prof.~\textsc{Encke} aus seinen eignen und den Mannheimer Beobachtungen
herausgebracht, und Hrn.~Hofr.~\textsc{Gauss} in einem Schreiben vom 5.~Febr. mitgetheilt hat, verdienen die
g"rosste Aufmerksamkeit. Obgleich diese Beobachtungen nur einen Zeitraum von drei Wochen umfassen,
so liessen sie sich doch nicht in einer Parabel ohne Fehler von mehreren Minuten darstellen, Fehler, die
bei diesen Beobachtungen nicht zul"assig waren. Nach mehreren Versuchen fand Hr.~\textsc{Encke} unter
Zuziehung der freilich nur sehr rohen \textsc{Pons}schen Angaben, eine \textit{elliptische Bahn von 3 Jahren und
7 Monaten} Umlaufszeit, welche alle Beobachtungen auf das trefflichste vereinigte.

Folgendes ist die Uebersicht der Unterschiede der in Mannheim und Seeberg beobachteten Oerter
von den berechneten:

\begin{center}
\small
\begin{tabular}{c|c|c|c|c}
\toprule
& & Ger.~Aufst. & Abweich. & Beobachter \\
\midrule
Dec. 22 & & $-\,1.0$ & $-\,5.5$ & \textsc{Nicolai} \\
\phantom{22} 25 & & $+\,7.6$ & $+17.3$ & \\
\phantom{22} 27 & & & $-\,\phantom{0}4.7$ & \\
\phantom{22} 29 & & & $-\,\phantom{0}1.8$ & \\
\phantom{22} & & $+\,0.3$ & & \\
Jan. \phantom{0}1 & & $+11.7$ & & \\
\phantom{0}\phantom{0} \phantom{0}4 & & $+14.8$ & & \textsc{Encke} \\
\phantom{0}\phantom{0} \phantom{0}5 & & $+11.2$ & & \\
\phantom{0}\phantom{0} \phantom{0}6 & & $+14.6$ & & \\
\phantom{0}\phantom{0} 12 & & $+26.5$ & & \\
\phantom{0}\phantom{0} & & $+15.6$ & & \\
\phantom{0}\phantom{0} & & & $+12.9$ & \\
\phantom{0}\phantom{0} & & & $+27.2$ & \\
\phantom{0}\phantom{0} & & $+23.8$ & & \\
\phantom{0}\phantom{0} & & $+30.6$ & & \\
\phantom{0}\phantom{0} & & $+36.1$ & & \\
\bottomrule
\end{tabular}
\end{center}

Ohne auf die \textsc{Pons}schen Angaben R"ucksicht zu nehmen, w"urde sich eine noch merklich bessere
Uebereinstimmung haben erreichen lassen.

Diese elliptischen Elemente sind folgende:

\begin{center}
\begin{tabular}{l@{\dotfill}l}
Zeit der Sonnenn"ahe 1819.~Januar & $27.13417^{*})$ \\
L"ange der Sonnenn"ahe & $156^{\circ}\;14'\;\phantom{0}8''$ \\
L"ange des Knoten & $334\;18\;\phantom{0}8$ \\
(beide auf das mittlere Aequinoctium von 1819 bezogen) & \\
Neigung der Bahn & $13^{\circ}\;42'\;30''$ \\
Logarithm der halben grossen Axe & $0.3697758$ \\
Excentricit"atswinkel & $\phantom{0}\phantom{0}29\;\phantom{0}1\;\phantom{0}\phantom{0}\phantom{0}$ \\
\end{tabular}
\end{center}

$^{*})$ wahrscheinlich Seeberger Zeit, was von Hrn.~\textsc{Encke} nicht ausdr"ucklich bemerkt war.

Hoffentlich werden in Zukunft noch fr"uhere gute Beobachtungen aus Frankreich oder Italien bekannt
werden, wodurch die Ungewissheit der Bahnbestimmung in engere Grenzen gebracht werden wird.
Noch wichtiger w"are es, wenn noch einige gute Beobachtungen nach dem Durchgange durch die Sonnenn"ahe
in s"udlichen Gegenden gemacht werden k"onnten. In unsern Breiten, wo der Comet sich nur noch
wenig "uber den s"udlichen Horizont erheben wird, ist dazu keine Hoffnung.

H"ochst merkw"urdig ist die grosse Uebereinstimmung dieser Elemente mit denen des ersten Cometen
vom Jahr 1805. Zur Vergleichung setzen wir dieselben, so wie sie damals von Hrn.~Prof.~\textsc{Bessel}
in einer Parabel berechnet wurden, hier her, und bemerken nur noch, dass schon damals die Vergleichung
der Beobachtungen mit der Rechnung fast unverkennbare Spuren einer Abweichung der Bahn jenes Cometen
von der Parabel zeigte, die jedoch von keinem Astronomen weiter verfolgt zu sein scheinen.

\begin{center}
\begin{tabular}{l@{\dotfill}l}
Durchgang durch die Sonnenn"ahe 1818 November 18.13782 & Pariser Zeit \\
L"ange der Sonnenn"ahe & $147^{\circ}\;51'\;28''$ \\
L"ange des aufsteigenden Knoten & $344\;37\;19$ \\
Neigung der Bahn & $\phantom{0}15\;36\;36$ \\
Logarithm des kleinsten Abstandes & $9.57820$ \\
\end{tabular}
\end{center}

(Bei den \textsc{Encke}schen Elementen des Cometen von 1818--1819 wird dieser Logarithm $9.52579$.)
Gewiss verdient die Frage, ob der jetzige Comet seit dem Jahr 1805 solche St"orungen erlitten haben kann,
dass die Unterschiede in den Elementen sich dadurch erkl"aren lassen, und ob demnach beide Cometen
f"ur identisch zu halten sind, eine sorgf"altige Untersuchung.

Den andern von \textsc{Pons} entdeckten Cometen fand seinerseits auch Hr.~Prof.~\textsc{Bessel} am 22.~December
auf der K"onigsberger Sternwarte auf.

Folgende Beobachtungen theilte dieser Astronom in einem Schreiben vom 14.~Januar dem Hrn.~Hofr.~\textsc{Gauss} mit:

\begin{center}
\small
\begin{tabular}{c|c|c|c|c}
\toprule
& Mittlere Zeit & & Gerade & Nordliche \\
& in K"onigsberg & & Aufsteigung & Abweichung \\
\midrule
1818 Dec. 22 & $5^{\mathrm{h}}\;36^{\mathrm{m}}\;51^{\mathrm{s}}$ & & $303^{\circ}\;10'\;21''$ & $36^{\circ}\;48'\;29''$ \\
\phantom{22} 24 & $6\;21\;35$ & & $303\;37\;29$ & $36\;51\;\phantom{0}\phantom{0}3$ \\
\phantom{22} & & & & \phantom{0}\phantom{0}\phantom{0}\phantom{0}$20.28$ \\
\phantom{22} 25 & $6\;22\;20$ & & $311\;56\;29$ & \phantom{0}\phantom{0}$\phantom{0}\phantom{0}48.29$ \\
\phantom{22} 26 & $6\;\phantom{0}4\;56$ & & $313\;17\;17.2$ & \phantom{0}$\phantom{0}37\;\phantom{0}7\;53.1$ \\
\phantom{22} & $6\;55\;\phantom{0}0$ & & $315\;38\;48.2$ & \phantom{0}$\phantom{0}37\;\phantom{0}4\;\phantom{0}1.2$ \\
\phantom{22} 28 & $6\;\phantom{0}0\;22$ & & $316\;38\;18.1$ & $36\;58\;35.2$ \\
\phantom{22} & & & & \phantom{0}\phantom{0}\phantom{0}\phantom{0}$45.6$ \\
1819 Jan. \phantom{0}2 & $6\;11\;11$ & & $319\;39\;24.48$ & $36\;53\;31.6$ \\
\phantom{0}\phantom{0} & & & $324\;39\;12.3$ & $36\;22\;34.2$ \\
\phantom{0}\phantom{0} & & & & \phantom{0}$\phantom{0}36\;15\;54.3$ \\
\phantom{0}\phantom{0} & & & $325\;22\;\phantom{0}3.4$ & \\
\bottomrule
\end{tabular}
\end{center}

\newpage

Folgende von Hrn.~Prof.~\textsc{Bessel} auf die Beobachtungen vom 22.~und 27.~December und vom
2.~Januar gegr"undete vorl"aufige parabolische Elemente waren diesem Schreiben beigef"ugt:

\begin{center}
\begin{tabular}{l@{\dotfill}l}
Zeit der Sonnenn"ahe 1818 December 4.0968 & Pariser Zeit. \\
Aufsteigender Knoten & $\phantom{0}9^{\circ}\;\phantom{0}7'$ \\
Neigung der Bahn & $117\;19$ \\
Abstand des Perihel vom Knoten & $347\;\phantom{0}0$ \\
Logarithm des kleinsten Abstandes & $9.928324$ \\
\end{tabular}
\end{center}

Nach Anleitung dieser Elemente wurde dieser Comet vom Hrn.~Hofr.~\textsc{Gauss} sogleich als ein sehr
blasser Nebelfleck aufgefunden; die Elemente gaben die gerade Aufsteigung gegen $40'$, die Abweichung
aber $10'$ zu gross. Hierauf wurden durch Hrn.~Prof.~\textsc{Harding} noch folgende Beobachtungen auf der
hiesigen Sternwarte angestellt:

\begin{center}
\small
\begin{tabular}{c|c|c|c|c}
\toprule
1819 & Mittlere Zeit & & Gerade & Nordliche \\
& in G"ottingen & & Aufsteigung & Abweichung \\
\midrule
Jan. 26 & $7^{\mathrm{h}}\;19^{\mathrm{m}}\;\phantom{0}5^{\mathrm{s}}$ & & $335^{\circ}\;36'\;47''$ & $35^{\circ}\;\phantom{0}\phantom{0}'\;29''$ \\
\phantom{26} 28 & $7\;\phantom{0}2\;35$ & & $335\;50\;38$ & $35\;18\;50$ \\
\phantom{26} & $7\;22\;16$ & & $336\;\phantom{0}4\;38$ & $35\;20\;58$ \\
\phantom{26} & $7\;27\;56$ & & $336\;16\;33$ & $35\;22\;32$ \\
\phantom{26} & $6\;54\;10$ & & & $35\;23\;31$ \\
\bottomrule
\end{tabular}
\end{center}

Bei der immer zunehmenden Lichtschw"ache dieses durch seine starke geocentrische Bewegung seit
der ersten Entdeckung merkw"urdigen Cometen werden schwerlich noch sp"atere Beobachtungen zu erwarten
sein. Ob auch bei diesem Cometen die vorhandenen Beobachtungen zur Bestimmung der Abweichung der
Bahn von der Parabel hinreichend, und welches die Resultate sein werden, wird eine k"unftige Untersuchung
lehren.

\newpage

\section*{COMET.}
\markboth{COMET.}{}

\begin{center}
\rule{0.6\textwidth}{0.4pt}\\[0.5em]
\textit{G"ottingische gelehrte Anzeigen. St"uck 83. Seite 825\,..\,829. 1819 Mai 24.}\\[0.3em]
\rule{0.6\textwidth}{0.4pt}
\end{center}

Im 28.~St"uck dieser Anzeigen [S.~417 d.~B.] haben wir eine Nachricht "uber die beiden Cometen
vom vorigen Jahre mitgetheilt, und dabei bereits die Wahrscheinlichkeit der Identit"at des einen dieses
Cometen mit dem ersten von 1805 angedeutet. Diese Wahrscheinlichkeit ist durch die sp"atern Rechnungen
des Hrn.~Prof.~\textsc{Encke} so sehr erh"ohet worden, dass kaum noch ein Zweifel an der Identit"at statt finden
kann. Wir tragen kein Bedenken, diese Bereicherung der Kenntniss unsers Sonnensystems zu den
merkw"urdigsten Entdeckungen des neunzehnten Jahrhunderts zu z"ahlen, und theilen hier dar"uber das
Haupts"achlichste, aus einigen Briefen jenes Astronomen an Hrn~Hofr.~\textsc{Gauss}, mit. Da noch immer
von dem Cometen von 1818 keine weitern brauchbaren Beobachtungen, als die im angef"uhrten Blatte
abgedruckten, bekannt geworden sind, so hielt Hr.~Prof.~\textsc{Encke} vorerst f"ur das Rathsamste, eine neue
Discussion der Beobachtungen des Cometen von 1805 vorzunehmen. Es zeigt sich hier das "uberraschende
Resultat, dass eine Ellipse mit einer Umlaufszeit von 3.42 Jahren die Beobachtungen am besten darstelle;
f"ur den Cometen von 1818 hatte er nach den angef"uhrten Elementen eine Umlaufszeit von 3.59 Jahren
gefunden: beide Angaben involvirten nat"urlich noch einige Unzuverl"assigkeit. Da nun zwischen den
beiden beobachteten Durchgangszeiten durch die Sonnenn"ahe 13.19 Jahre verflossen sind, so war unter der
Voraussetzung der Identit"at das Wahrscheinlichste, dass der Comet in der Zwischenzeit vier Uml"aufe
gemacht habe. (Dass der, wenn auch nicht zu den allerschw"achsten geh"orende, doch nur telescopische
Comet, bei seiner dreimaligen Wiederkunft zur Sonnenn"ahe 1809, 1812 und 1815 unbemerkt geblieben
ist, kann keinen Unterrichteten befremden). Mit einer diesem Schl"usse gem"ass abge"anderten Umlaufszeit
berechnete Hr.~Prof.~\textsc{Encke} neue elliptische Elemente nach den Beobachtungen von 1805, wodurch
diese auf das befriedigendste dargestellt wurden. Ueberraschend angenehm war hiebei, dass diese neue
Bahn in den einzelnen Elementen der aus den Beobachtungen von 1818 berechneten ungemein viel
n"aher gekommen war. Hr.~\textsc{Encke} entschloss sich, nach diesem aufmunternden Erfolge, die St"orungen,
welche jene Elemente w"ahrend der ganzen 13 Jahre durch die Einwirkung des Jupiter erleiden mussten,
nach der Methode des Hrn.~Hofrath \textsc{Gauss} zu berechnen. Auch diese m"uhsame Arbeit blieb nicht
unbelohnt: eine noch bessere Uebereinstimmung war die Folge davon. Wir setzen hier diese Elemente
her, wie sie also, blos auf die Beobachtungen von 1805 gegr"undet, und nur in R"ucksicht der grossen
Axe so bestimmt, dass nach vier Uml"aufen die Durchgangszeit durch die Sonnenn"ahe mit der 1819
beobachteten zusammenfiel, nach den Einwirkungen des Jupiter f"ur den Anfang des Jahrs 1819

\newpage

sich ergeben haben:

\begin{center}
\begin{tabular}{l@{\dotfill}l}
Zeit der Sonnenn"ahe & $1819.\ 182^{\mathrm{t}}\;15^{\mathrm{h}}\;21^{\mathrm{m}}$~M.~Z.~Berlin \\
L"ange der Sonnenn"ahe & $147^{\circ}\;52'\;31''$ \\
L"ange des aufsteigenden Knoten & $334\;35\;22$ \\
Neigung der Bahn & $13\;28\;53$ \\
Logarithm der halben grossen Axe & $0.369558$ \\
Excentricit"at & $0.84953$ \\
S"achliche Umlaufszeit & $1205^{\mathrm{t}}\;21^{\mathrm{h}}$ \\
\end{tabular}
\end{center}

Bei diesen Elementen sind die beobachteten und berechneten Oerter f"ur 1805 verglichen:

\begin{center}
\small
\begin{tabular}{c|c|c|c|c}
\toprule
& \multicolumn{2}{c|}{Ger.~Aufst.} & \multicolumn{2}{c}{Abweich.} \\
& beob. & ber. & beob. & ber. \\
\midrule
Aug. 19 & $197^{\circ}\;20'$ & $197^{\circ}\;23'$ & $16^{\circ}\;20'$~N. & $16^{\circ}\;15'$~N. \\
Sept. 2 & $202\;12$ & $202\;10$ & $11\;26$ & $11\;29$ \\
Oct. 1 & $209\;\phantom{0}8$ & $209\;10$ & $2\;50$ & $2\;50$ \\
Nov. 1 & $216\;12$ & $216\;11$ & $6\;11$~S. & $6\;12$ \\
Dec. 24 & $243\;29$ & $243\;29$ & $28\;38$ & $28\;39$ \\
\bottomrule
\end{tabular}
\end{center}

Hieraus sind die Unterschiede der Elemente des Cometen von 1805 und 1818 nach Abzug der
St"orungen:

\begin{center}
\begin{tabular}{l|c|c}
\toprule
& 1805 & 1818 \\
\midrule
L"ange der Sonnenn"ahe & $147^{\circ}\;52'$ & $147^{\circ}\;51'$ \\
L"ange des Knoten & $334\;35$ & $334\;37$ \\
Neigung & $13\;29$ & $15\;36$ \\
Log.~halbe grosse Axe & $0.36956$ & $0.36978$ \\
Excentricit"at & $0.8495$ & $0.8456$ \\
\bottomrule
\end{tabular}
\end{center}

Diese Elemente sind nun in der That ungemein nahe mit einander "ubereinstimmend; allein die
Neigung scheint eine zu betr"achtliche Ver"anderung erlitten zu haben, als dass sie durch die St"orungen
erkl"art werden k"onnte. Dieser Umstand verdient also noch eine weitere Untersuchung.

\newpage

\section*{PALLAS.}
\markboth{PALLAS.}{}
\addcontentsline{toc}{section}{Pallas}

\begin{center}
\rule{0.6\textwidth}{0.4pt}\\[0.5em]
\textit{G"ottingische gelehrte Anzeigen. St"uck 131. Seite 1305\,..\,1306. 1819 Sept.~13.}\\[0.3em]
\rule{0.6\textwidth}{0.4pt}
\end{center}

Bei der im vorigen Jahre bemerkten g"unstigen Opposition der \textsc{Pallas} ist dieselbe auf der hiesigen
Sternwarte mit dem \textsc{Repsold}schen Kreise t"aglich, bei sch"onem Wetter, beobachtet worden. Die
Beobachtungen sind theils absolut, theils differential mittelst der \textsc{Schumacher}schen \textit{Rectascensions- und
Declinations-Tafeln} f"ur den Meridiankreis, und die \textit{Instrumental-Constanten} reducirt. Hr.~Hofr.~\textsc{Gauss}
hat die hiesigen Beobachtungen mit den \textsc{Tittel}schen Elementen verglichen, und dabei zugleich die Oerter
f"ur einige andre Tage berechnet, wodurch die Opposition noch etwas genauer bestimmt wird. Das
Resultat ist folgendes:

\begin{center}
\begin{tabular}{l@{\quad}l}
1819 Sept.~8. & $18^{\mathrm{h}}\;44^{\mathrm{m}}\;17^{\mathrm{s}}$~M.~Z.~G"ottingen \\
Wahre L"ange \dotfill & $345^{\circ}\;55'\;42''$ \\
Geocentrische Breite \dotfill & $\phantom{0}6\;46\;57.5$ \\
\end{tabular}
\end{center}

Die Vergleichung mit den \textsc{Tittel}schen Elementen gibt die Unterschiede (Beob.~$-$~Rechn.):

\begin{center}
\small
\begin{tabular}{c|c|c}
\toprule
& Ger.~Aufst. & Abweich. \\
\midrule
Juni \phantom{0}3 & $-\,4''4$ & $+\,2''5$ \\
\phantom{0}\phantom{0} 11 & $-\,5.2$ & $+\,0.9$ \\
\phantom{0}\phantom{0} 17 & $-\,4.0$ & $+\,1.3$ \\
\phantom{0}\phantom{0} 25 & $-\,4.5$ & $+\,1.5$ \\
Juli \phantom{0}2 & $-\,5.3$ & $+\,0.8$ \\
\phantom{0}\phantom{0} 10 & $-\,4.8$ & $+\,1.0$ \\
\phantom{0}\phantom{0} 22 & $-\,4.2$ & $+\,1.2$ \\
Aug. \phantom{0}1 & $-\,4.6$ & $+\,1.1$ \\
\phantom{0}\phantom{0} 12 & $-\,4.3$ & $+\,0.7$ \\
Sept. \phantom{0}2 & $-\,4.0$ & $+\,0.5$ \\
\phantom{0}\phantom{0} \phantom{0}7 & $-\,3.8$ & $+\,0.6$ \\
\bottomrule
\end{tabular}
\end{center}

\vspace{1em}

\begin{center}
\rule{0.6\textwidth}{0.4pt}\\[0.5em]
\textit{G"ottingische gelehrte Anzeigen. St"uck 187. Seite 1865\,..\,1866. 1820.~November 20.}\\[0.3em]
\rule{0.6\textwidth}{0.4pt}
\end{center}

Se.~K"onigl.~Hoheit der Herzog von \textsc{Clarence} hat der Universit"at einen Beweis seines gn"adigen
Andenkens gegeben, indem Er die Bibliothek mit einer kostbaren Sammlung von Seecharten in 182
Bl"attern beschenkt hat. Es sind dies die in London in dem Hydrographical Office erscheinenden, und
mit dessen Stempel bezeichneten Charten; die nicht in den Buchhandel kommen, sondern nur f"ur den
Gebrauch der K"onigl.~Marine bestimmt sind. Die Sammlung umfasst nicht bloss die Europ"aischen
Gew"asser, sondern auch den gr"ossten Theil der K"usten von Africa, America, Ost- und Westindien, von
denen der H\kern-0.5pt{}r.~\textsc{Hydrograph}\kern-0.5pt{} der K"onigl.~Marine sie mittheilte.

\newpage

\vspace{2em}

\begin{center}
\rule{0.8\textwidth}{0.4pt}
\end{center}

\vspace{1em}

\textit{Ende des neunten Theils.}

\end{document}
